\pdfoutput=1
%% Author: PGL  Porta Mana
%% Created: 2015-05-01T20:53:34+0200
%% Last-Updated: 2020-02-03T11:24:40+0100
%%%%%%%%%%%%%%%%%%%%%%%%%%%%%%%%%%%%%%%%%%%%%%%%%%%%%%%%%%%%%%%%%%%%%%%%%%%%
\newif\ifarxiv
\arxivfalse
\ifarxiv\pdfmapfile{+classico.map}\fi
\newif\ifafour
\afourfalse% true = A4, false = A5
\newif\iftypodisclaim % typographical disclaim on the side
\typodisclaimtrue
\newcommand*{\memfontfamily}{zplx}
\newcommand*{\memfontpack}{newpxtext}
\documentclass[\ifafour a4paper,12pt,\else a5paper,10pt,\fi%extrafontsizes,%
onecolumn,oneside,article,%french,italian,german,swedish,latin,
british%
]{memoir}
\newcommand*{\updated}{\today}
\newcommand*{\firstdraft}{4 November 2015}
\newcommand*{\firstpublished}{4 November 2015}
\newcommand*{\propertitle}{Inferring the total activity of a large neuronal population\\ from a small sample}
\newcommand*{\pdftitle}{\propertitle}
\newcommand*{\headtitle}{Inferences about large neuronal populations}
\newcommand*{\pdfauthor}{P.G.L.  Porta Mana, Y. Roudi, E. Torre, V. Rostami}
\newcommand*{\headauthor}{Porta Mana \etal}
\newcommand*{\reporthead}{}% Report number

%%%%%%%%%%%%%%%%%%%%%%%%%%%%%%%%%%%%%%%%%%%%%%%%%%%%%%%%%%%%%%%%%%%%%%%%%%%%
%%% Calls to packages (uncomment as needed)
%%%%%%%%%%%%%%%%%%%%%%%%%%%%%%%%%%%%%%%%%%%%%%%%%%%%%%%%%%%%%%%%%%%%%%%%%%%%

%\usepackage{pifont}

%\usepackage{fontawesome}

\usepackage[T1]{fontenc} 
\input{glyphtounicode} \pdfgentounicode=1

\usepackage[utf8]{inputenx}

%\usepackage{newunicodechar}
% \newunicodechar{Ĕ}{\u{E}}
% \newunicodechar{ĕ}{\u{e}}
% \newunicodechar{Ĭ}{\u{I}}
% \newunicodechar{ĭ}{\u{\i}}
% \newunicodechar{Ŏ}{\u{O}}
% \newunicodechar{ŏ}{\u{o}}
% \newunicodechar{Ŭ}{\u{U}}
% \newunicodechar{ŭ}{\u{u}}
% \newunicodechar{Ā}{\=A}
% \newunicodechar{ā}{\=a}
% \newunicodechar{Ē}{\=E}
% \newunicodechar{ē}{\=e}
% \newunicodechar{Ī}{\=I}
% \newunicodechar{ī}{\={\i}}
% \newunicodechar{Ō}{\=O}
% \newunicodechar{ō}{\=o}
% \newunicodechar{Ū}{\=U}
% \newunicodechar{ū}{\=u}
% \newunicodechar{Ȳ}{\=Y}
% \newunicodechar{ȳ}{\=y}

\newcommand*{\bmmax}{0} % reduce number of bold fonts, before font packages
\newcommand*{\hmmax}{0} % reduce number of heavy fonts, before font packages

\usepackage{textcomp}

%\usepackage[normalem]{ulem}% package for underlining
% \makeatletter
% \def\ssout{\bgroup \ULdepth=-.35ex%\UL@setULdepth
%  \markoverwith{\lower\ULdepth\hbox
%    {\kern-.03em\vbox{\hrule width.2em\kern1.2\p@\hrule}\kern-.03em}}%
%  \ULon}
% \makeatother

\usepackage{amsmath}

\usepackage{mathtools}
\addtolength{\jot}{\jot} % increase spacing in multiline formulae
\setlength{\multlinegap}{0pt}

%\usepackage{empheq}% automatically calls amsmath and mathtools
%\newcommand*{\widefbox}[1]{\fbox{\hspace{1em}#1\hspace{1em}}}

%\usepackage{fancybox}

%\usepackage{framed}

% \usepackage[misc]{ifsym} % for dice
% \newcommand*{\diceone}{{\scriptsize\Cube{1}}}

\usepackage{amssymb}

\usepackage{amsxtra}

\usepackage[main=british,french,italian,german,swedish,latin,esperanto]{babel}\selectlanguage{british}
\newcommand*{\langfrench}{\foreignlanguage{french}}
\newcommand*{\langgerman}{\foreignlanguage{german}}
\newcommand*{\langitalian}{\foreignlanguage{italian}}
\newcommand*{\langswedish}{\foreignlanguage{swedish}}
\newcommand*{\langlatin}{\foreignlanguage{latin}}
\newcommand*{\langnohyph}{\foreignlanguage{nohyphenation}}

\usepackage[autostyle=false,autopunct=false,english=british]{csquotes}
\setquotestyle{british}

\usepackage{amsthm}
\newcommand*{\QED}{\textsc{q.e.d.}}
\renewcommand*{\qedsymbol}{\QED}
\theoremstyle{remark}
\newtheorem{note}{Note}
\newtheorem*{remark}{Note}
\newtheoremstyle{innote}{\parsep}{\parsep}{\footnotesize}{}{}{}{0pt}{}
\theoremstyle{innote}
\newtheorem*{innote}{}

\usepackage[shortlabels,inline]{enumitem}
\SetEnumitemKey{para}{itemindent=\parindent,leftmargin=0pt,listparindent=\parindent,parsep=0pt,itemsep=\topsep}
% \begin{asparaenum} = \begin{enumerate}[para]
% \begin{inparaenum} = \begin{enumerate*}
\setlist[enumerate,2]{label=\alph*.}
\setlist[enumerate]{label=\arabic*.,leftmargin=1.5\parindent}
\setlist[itemize]{leftmargin=1.5\parindent}
\setlist[description]{leftmargin=1.5\parindent}
% old alternative:
% \setlist[enumerate,2]{label=\alph*.}
% \setlist[enumerate]{leftmargin=\parindent}
% \setlist[itemize]{leftmargin=\parindent}
% \setlist[description]{leftmargin=\parindent}

\usepackage[babel,theoremfont,largesc]{newpxtext}

\usepackage[bigdelims,nosymbolsc%,smallerops % probably arXiv doesn't have it
]{newpxmath}
\linespread{1.083}%\useosf
%% smaller operators for old version of newpxmath
\makeatletter
\def\re@DeclareMathSymbol#1#2#3#4{%
    \let#1=\undefined
    \DeclareMathSymbol{#1}{#2}{#3}{#4}}
%\re@DeclareMathSymbol{\bigsqcupop}{\mathop}{largesymbols}{"46}
%\re@DeclareMathSymbol{\bigodotop}{\mathop}{largesymbols}{"4A}
\re@DeclareMathSymbol{\bigoplusop}{\mathop}{largesymbols}{"4C}
\re@DeclareMathSymbol{\bigotimesop}{\mathop}{largesymbols}{"4E}
\re@DeclareMathSymbol{\sumop}{\mathop}{largesymbols}{"50}
\re@DeclareMathSymbol{\prodop}{\mathop}{largesymbols}{"51}
\re@DeclareMathSymbol{\bigcupop}{\mathop}{largesymbols}{"53}
\re@DeclareMathSymbol{\bigcapop}{\mathop}{largesymbols}{"54}
%\re@DeclareMathSymbol{\biguplusop}{\mathop}{largesymbols}{"55}
\re@DeclareMathSymbol{\bigwedgeop}{\mathop}{largesymbols}{"56}
\re@DeclareMathSymbol{\bigveeop}{\mathop}{largesymbols}{"57}
%\re@DeclareMathSymbol{\bigcupdotop}{\mathop}{largesymbols}{"DF}
%\re@DeclareMathSymbol{\bigcapplusop}{\mathop}{largesymbolsPXA}{"00}
%\re@DeclareMathSymbol{\bigsqcupplusop}{\mathop}{largesymbolsPXA}{"02}
%\re@DeclareMathSymbol{\bigsqcapplusop}{\mathop}{largesymbolsPXA}{"04}
%\re@DeclareMathSymbol{\bigsqcapop}{\mathop}{largesymbolsPXA}{"06}
\re@DeclareMathSymbol{\bigtimesop}{\mathop}{largesymbolsPXA}{"10}
%\re@DeclareMathSymbol{\coprodop}{\mathop}{largesymbols}{"60}
%\re@DeclareMathSymbol{\varprod}{\mathop}{largesymbolsPXA}{16}
\makeatother
%%
%% With euler font cursive for Greek letters - the [1] means 100% scaling
\DeclareFontFamily{U}{egreek}{\skewchar\font'177}%
\DeclareFontShape{U}{egreek}{m}{n}{<-6>s*[1]eurm5 <6-8>s*[1]eurm7 <8->s*[1]eurm10}{}%
\DeclareFontShape{U}{egreek}{m}{it}{<->s*[1]eurmo10}{}%
\DeclareFontShape{U}{egreek}{b}{n}{<-6>s*[1]eurb5 <6-8>s*[1]eurb7 <8->s*[1]eurb10}{}%
\DeclareFontShape{U}{egreek}{b}{it}{<->s*[1]eurbo10}{}%
\DeclareSymbolFont{egreeki}{U}{egreek}{m}{it}%
\SetSymbolFont{egreeki}{bold}{U}{egreek}{b}{it}% from the amsfonts package
\DeclareSymbolFont{egreekr}{U}{egreek}{m}{n}%
\SetSymbolFont{egreekr}{bold}{U}{egreek}{b}{n}% from the amsfonts package
% Take also \sum, \prod, \coprod symbols from Euler fonts
\DeclareFontFamily{U}{egreekx}{\skewchar\font'177}
\DeclareFontShape{U}{egreekx}{m}{n}{%
       <-7.5>s*[0.9]euex7%
    <7.5-8.5>s*[0.9]euex8%
    <8.5-9.5>s*[0.9]euex9%
    <9.5->s*[0.9]euex10%
}{}
\DeclareSymbolFont{egreekx}{U}{egreekx}{m}{n}
\DeclareMathSymbol{\sumop}{\mathop}{egreekx}{"50}
\DeclareMathSymbol{\prodop}{\mathop}{egreekx}{"51}
\DeclareMathSymbol{\coprodop}{\mathop}{egreekx}{"60}
\makeatletter
\def\sum{\DOTSI\sumop\slimits@}
\def\prod{\DOTSI\prodop\slimits@}
\def\coprod{\DOTSI\coprodop\slimits@}
\makeatother
% Greek letters not usually given in LaTeX.
\DeclareMathSymbol{\varpartial}{\mathalpha}{egreeki}{"40}
\DeclareMathSymbol{\partialup}{\mathalpha}{egreekr}{"40}
\DeclareMathSymbol{\alpha}{\mathalpha}{egreeki}{"0B}
\DeclareMathSymbol{\beta}{\mathalpha}{egreeki}{"0C}
\DeclareMathSymbol{\gamma}{\mathalpha}{egreeki}{"0D}
\DeclareMathSymbol{\delta}{\mathalpha}{egreeki}{"0E}
\DeclareMathSymbol{\epsilon}{\mathalpha}{egreeki}{"0F}
\DeclareMathSymbol{\zeta}{\mathalpha}{egreeki}{"10}
\DeclareMathSymbol{\eta}{\mathalpha}{egreeki}{"11}
\DeclareMathSymbol{\theta}{\mathalpha}{egreeki}{"12}
\DeclareMathSymbol{\iota}{\mathalpha}{egreeki}{"13}
\DeclareMathSymbol{\kappa}{\mathalpha}{egreeki}{"14}
\DeclareMathSymbol{\lambda}{\mathalpha}{egreeki}{"15}
\DeclareMathSymbol{\mu}{\mathalpha}{egreeki}{"16}
\DeclareMathSymbol{\nu}{\mathalpha}{egreeki}{"17}
\DeclareMathSymbol{\xi}{\mathalpha}{egreeki}{"18}
\DeclareMathSymbol{\omicron}{\mathalpha}{egreeki}{"6F}
\DeclareMathSymbol{\pi}{\mathalpha}{egreeki}{"19}
\DeclareMathSymbol{\rho}{\mathalpha}{egreeki}{"1A}
\DeclareMathSymbol{\sigma}{\mathalpha}{egreeki}{"1B}
 \DeclareMathSymbol{\tau}{\mathalpha}{egreeki}{"1C}
\DeclareMathSymbol{\upsilon}{\mathalpha}{egreeki}{"1D}
\DeclareMathSymbol{\phi}{\mathalpha}{egreeki}{"1E}
\DeclareMathSymbol{\chi}{\mathalpha}{egreeki}{"1F}
\DeclareMathSymbol{\psi}{\mathalpha}{egreeki}{"20}
\DeclareMathSymbol{\omega}{\mathalpha}{egreeki}{"21}
\DeclareMathSymbol{\varepsilon}{\mathalpha}{egreeki}{"22}
\DeclareMathSymbol{\vartheta}{\mathalpha}{egreeki}{"23}
\DeclareMathSymbol{\varpi}{\mathalpha}{egreeki}{"24}
\let\varrho\rho 
\let\varsigma\sigma
 \let\varkappa\kappa
\DeclareMathSymbol{\varphi}{\mathalpha}{egreeki}{"27}
%
\DeclareMathSymbol{\varAlpha}{\mathalpha}{egreeki}{"41}
\DeclareMathSymbol{\varBeta}{\mathalpha}{egreeki}{"42}
\DeclareMathSymbol{\varGamma}{\mathalpha}{egreeki}{"00}
\DeclareMathSymbol{\varDelta}{\mathalpha}{egreeki}{"01}
\DeclareMathSymbol{\varEpsilon}{\mathalpha}{egreeki}{"45}
\DeclareMathSymbol{\varZeta}{\mathalpha}{egreeki}{"5A}
\DeclareMathSymbol{\varEta}{\mathalpha}{egreeki}{"48}
\DeclareMathSymbol{\varTheta}{\mathalpha}{egreeki}{"02}
 \DeclareMathSymbol{\varIota}{\mathalpha}{egreeki}{"49}
\DeclareMathSymbol{\varKappa}{\mathalpha}{egreeki}{"4B}
\DeclareMathSymbol{\varLambda}{\mathalpha}{egreeki}{"03}
\DeclareMathSymbol{\varMu}{\mathalpha}{egreeki}{"4D}
\DeclareMathSymbol{\varNu}{\mathalpha}{egreeki}{"4E}
\DeclareMathSymbol{\varXi}{\mathalpha}{egreeki}{"04}
\DeclareMathSymbol{\varOmicron}{\mathalpha}{egreeki}{"4F}
\DeclareMathSymbol{\varPi}{\mathalpha}{egreeki}{"05}
\DeclareMathSymbol{\varRho}{\mathalpha}{egreeki}{"50}
\DeclareMathSymbol{\varSigma}{\mathalpha}{egreeki}{"06}
\DeclareMathSymbol{\varTau}{\mathalpha}{egreeki}{"54}
\DeclareMathSymbol{\varUpsilon}{\mathalpha}{egreeki}{"07}
\DeclareMathSymbol{\varPhi}{\mathalpha}{egreeki}{"08}
\DeclareMathSymbol{\varChi}{\mathalpha}{egreeki}{"58}
\DeclareMathSymbol{\varPsi}{\mathalpha}{egreeki}{"09}
\DeclareMathSymbol{\varOmega}{\mathalpha}{egreeki}{"0A} 
%
\DeclareMathSymbol{\Alpha}{\mathalpha}{egreekr}{"41}
\DeclareMathSymbol{\Beta}{\mathalpha}{egreekr}{"42}
\DeclareMathSymbol{\Gamma}{\mathalpha}{egreekr}{"00}
\DeclareMathSymbol{\Delta}{\mathalpha}{egreekr}{"01}
\DeclareMathSymbol{\Epsilon}{\mathalpha}{egreekr}{"45}
\DeclareMathSymbol{\Zeta}{\mathalpha}{egreekr}{"5A}
\DeclareMathSymbol{\Eta}{\mathalpha}{egreekr}{"48}
\DeclareMathSymbol{\Theta}{\mathalpha}{egreekr}{"02}
\DeclareMathSymbol{\Iota}{\mathalpha}{egreekr}{"49}
\DeclareMathSymbol{\Kappa}{\mathalpha}{egreekr}{"4B}
\DeclareMathSymbol{\Lambda}{\mathalpha}{egreekr}{"03}
\DeclareMathSymbol{\Mu}{\mathalpha}{egreekr}{"4D}
\DeclareMathSymbol{\Nu}{\mathalpha}{egreekr}{"4E}
\DeclareMathSymbol{\Xi}{\mathalpha}{egreekr}{"04}
\DeclareMathSymbol{\Omicron}{\mathalpha}{egreekr}{"4F}
\DeclareMathSymbol{\Pi}{\mathalpha}{egreekr}{"05}
\DeclareMathSymbol{\Rho}{\mathalpha}{egreekr}{"50}
\DeclareMathSymbol{\Sigma}{\mathalpha}{egreekr}{"06}
\DeclareMathSymbol{\Tau}{\mathalpha}{egreekr}{"54}
\DeclareMathSymbol{\Upsilon}{\mathalpha}{egreekr}{"07}
\DeclareMathSymbol{\Phi}{\mathalpha}{egreekr}{"08}
\DeclareMathSymbol{\Chi}{\mathalpha}{egreekr}{"58}
\DeclareMathSymbol{\Psi}{\mathalpha}{egreekr}{"09}
\DeclareMathSymbol{\Omega}{\mathalpha}{egreekr}{"0A}
%
\DeclareMathSymbol{\alphaup}{\mathalpha}{egreekr}{"0B}
\DeclareMathSymbol{\betaup}{\mathalpha}{egreekr}{"0C}
\DeclareMathSymbol{\gammaup}{\mathalpha}{egreekr}{"0D}
 \DeclareMathSymbol{\deltaup}{\mathalpha}{egreekr}{"0E}
\DeclareMathSymbol{\epsilonup}{\mathalpha}{egreekr}{"0F}
\DeclareMathSymbol{\zetaup}{\mathalpha}{egreekr}{"10}
\DeclareMathSymbol{\etaup}{\mathalpha}{egreekr}{"11}
\DeclareMathSymbol{\thetaup}{\mathalpha}{egreekr}{"12}
\DeclareMathSymbol{\iotaup}{\mathalpha}{egreekr}{"13}
\DeclareMathSymbol{\kappaup}{\mathalpha}{egreekr}{"14}
\DeclareMathSymbol{\lambdaup}{\mathalpha}{egreekr}{"15}
\DeclareMathSymbol{\muup}{\mathalpha}{egreekr}{"16}
\DeclareMathSymbol{\nuup}{\mathalpha}{egreekr}{"17}
\DeclareMathSymbol{\xiup}{\mathalpha}{egreekr}{"18}
\DeclareMathSymbol{\omicronup}{\mathalpha}{egreekr}{"6F}
  \DeclareMathSymbol{\piup}{\mathalpha}{egreekr}{"19}
\DeclareMathSymbol{\rhoup}{\mathalpha}{egreekr}{"1A}
\DeclareMathSymbol{\sigmaup}{\mathalpha}{egreekr}{"1B}
\DeclareMathSymbol{\tauup}{\mathalpha}{egreekr}{"1C}
\DeclareMathSymbol{\upsilonup}{\mathalpha}{egreekr}{"1D}
\DeclareMathSymbol{\phiup}{\mathalpha}{egreekr}{"1E}
\DeclareMathSymbol{\chiup}{\mathalpha}{egreekr}{"1F}
\DeclareMathSymbol{\psiup}{\mathalpha}{egreekr}{"20}
\DeclareMathSymbol{\omegaup}{\mathalpha}{egreekr}{"21}
\DeclareMathSymbol{\varepsilonup}{\mathalpha}{egreekr}{"22}
\DeclareMathSymbol{\varthetaup}{\mathalpha}{egreekr}{"23}
\DeclareMathSymbol{\varpiup}{\mathalpha}{egreekr}{"24}
\let\varrhoup\rhoup 
\let\varsigmaup\sigmaup
\let\varkappaup\kappaup
\DeclareMathSymbol{\varphiup}{\mathalpha}{egreekr}{"27}
% Greek letters not usually given in LaTeX.

%\usepackage%[scaled=0.9]%
%{classico}%  Optima as sans-serif font
\renewcommand\sfdefault{uop}
\DeclareMathAlphabet{\mathsf}  {T1}{\sfdefault}{m}{sl}
\SetMathAlphabet{\mathsf}{bold}{T1}{\sfdefault}{b}{sl}
\newcommand*{\mathte}[1]{\textbf{\textit{\textsf{#1}}}}
% Upright sans-serif math alphabet
% \DeclareMathAlphabet{\mathsu}  {T1}{\sfdefault}{m}{n}
% \SetMathAlphabet{\mathsu}{bold}{T1}{\sfdefault}{b}{n}

% DejaVu Mono as typewriter text
\usepackage[scaled=0.84]{DejaVuSansMono}

\usepackage{mathdots}

\usepackage[usenames]{xcolor}
% Tol (2012) colour-blind-, print-, screen-friendly colours, alternative scheme; Munsell terminology
\definecolor{mypurpleblue}{RGB}{68,119,170}
\definecolor{myblue}{RGB}{102,204,238}
\definecolor{mygreen}{RGB}{34,136,51}
\definecolor{myyellow}{RGB}{204,187,68}
\definecolor{myred}{RGB}{238,102,119}
\definecolor{myredpurple}{RGB}{170,51,119}
\definecolor{mygrey}{RGB}{187,187,187}
% Tol (2012) colour-blind-, print-, screen-friendly colours; Munsell terminology
% \definecolor{lbpurple}{RGB}{51,34,136}
% \definecolor{lblue}{RGB}{136,204,238}
% \definecolor{lbgreen}{RGB}{68,170,153}
% \definecolor{lgreen}{RGB}{17,119,51}
% \definecolor{lgyellow}{RGB}{153,153,51}
% \definecolor{lyellow}{RGB}{221,204,119}
% \definecolor{lred}{RGB}{204,102,119}
% \definecolor{lpred}{RGB}{136,34,85}
% \definecolor{lrpurple}{RGB}{170,68,153}
\definecolor{lgrey}{RGB}{221,221,221}
%\newcommand*\mycolourbox[1]{%
%\colorbox{mygrey}{\hspace{1em}#1\hspace{1em}}}
\colorlet{shadecolor}{lgrey}

\usepackage{bm}

\usepackage{microtype}

\usepackage[backend=biber,mcite,%subentry,
citestyle=authoryear-comp,bibstyle=pglpm-authoryear,autopunct=false,sorting=ny,sortcites=false,natbib=false,maxcitenames=1,maxbibnames=8,minbibnames=8,giveninits=true,uniquename=false,uniquelist=false,maxalphanames=1,block=space,hyperref=true,defernumbers=false,useprefix=true,sortupper=false,language=british,parentracker=false]{biblatex}
\DeclareSortingScheme{ny}{\sort{\field{sortname}\field{author}\field{editor}}\sort{\field{year}}}
\iffalse\makeatletter%%% replace parenthesis with brackets
\newrobustcmd*{\parentexttrack}[1]{%
  \begingroup
  \blx@blxinit
  \blx@setsfcodes
  \blx@bibopenparen#1\blx@bibcloseparen
  \endgroup}
\AtEveryCite{%
  \let\parentext=\parentexttrack%
  \let\bibopenparen=\bibopenbracket%
  \let\bibcloseparen=\bibclosebracket}
\makeatother\fi
\DefineBibliographyExtras{british}{\def\finalandcomma{\addcomma}}
\renewcommand*{\finalnamedelim}{\addcomma\space}
\setcounter{biburlnumpenalty}{1}
\setcounter{biburlucpenalty}{0}
\setcounter{biburllcpenalty}{1}
\DeclareDelimFormat{multicitedelim}{\addsemicolon\space}
\DeclareDelimFormat{compcitedelim}{\addsemicolon\space}
\DeclareDelimFormat{postnotedelim}{\space}
\ifarxiv\else\addbibresource{portamanabib.bib}\fi
\renewcommand{\bibfont}{\footnotesize}
%\appto{\citesetup}{\footnotesize}% smaller font for citations
\defbibheading{bibliography}[\bibname]{\section*{#1}\addcontentsline{toc}{section}{#1}%\markboth{#1}{#1}
}
\newcommand*{\citep}{\footcites}
\newcommand*{\citey}{\parencites*}
%\renewcommand*{\cite}{\parencite}
%\newcommand*{\citeps}{\parencites}
\providecommand{\href}[2]{#2}
\providecommand{\eprint}[2]{\texttt{\href{#1}{#2}}}
\newcommand*{\amp}{\&}
% \newcommand*{\citein}[2][]{\textnormal{\textcite[#1]{#2}}%\addtocategory{extras}{#2}
% }
\newcommand*{\citein}[2][]{\textnormal{\textcite[#1]{#2}}%\addtocategory{extras}{#2}
}
\newcommand*{\citebi}[2][]{\textcite[#1]{#2}%\addtocategory{extras}{#2}
}
\newcommand*{\subtitleproc}[1]{}
\newcommand*{\chapb}{ch.}
%
% \def\arxivp{}
% \def\mparcp{}
% \def\philscip{}
% \def\biorxivp{}
% \newcommand*{\arxivsi}{\texttt{arXiv} eprints available at \url{http://arxiv.org/}.\\}
% \newcommand*{\mparcsi}{\texttt{mp\_arc} eprints available at \url{http://www.ma.utexas.edu/mp_arc/}.\\}
% \newcommand*{\philscisi}{\texttt{philsci} eprints available at \url{http://philsci-archive.pitt.edu/}.\\}
% \newcommand*{\biorxivsi}{\texttt{bioRxiv} eprints available at \url{http://biorxiv.org/}.\\}
\newcommand*{\arxiveprint}[1]{%\global\def\arxivp{\arxivsi}%\citeauthor{0arxivcite}\addtocategory{ifarchcit}{0arxivcite}%eprint
\texttt{arXiv:\urlalt{https://arxiv.org/abs/#1}{#1}}%
%\texttt{\href{http://arxiv.org/abs/#1}{\protect\url{arXiv:#1}}}%
%\renewcommand{\arxivnote}{\texttt{arXiv} eprints available at \url{http://arxiv.org/}.}
}
\newcommand*{\mparceprint}[1]{%\global\def\mparcp{\mparcsi}%\citeauthor{0mparccite}\addtocategory{ifarchcit}{0mparccite}%eprint
\texttt{mp\_arc:\urlalt{http://www.ma.utexas.edu/mp_arc-bin/mpa?yn=#1}{#1}}%
%\texttt{\href{http://www.ma.utexas.edu/mp_arc-bin/mpa?yn=#1}{\protect\url{mp_arc:#1}}}%
%\providecommand{\mparcnote}{\texttt{mp_arc} eprints available at \url{http://www.ma.utexas.edu/mp_arc/}.}
}
\newcommand*{\haleprint}[1]{%\global\def\arxivp{\arxivsi}%\citeauthor{0arxivcite}\addtocategory{ifarchcit}{0arxivcite}%eprint
\texttt{HAL:\urlalt{https://hal.archives-ouvertes.fr/#1}{#1}}%
%\texttt{\href{http://arxiv.org/abs/#1}{\protect\url{arXiv:#1}}}%
%\renewcommand{\arxivnote}{\texttt{arXiv} eprints available at \url{http://arxiv.org/}.}
}
\newcommand*{\philscieprint}[1]{%\global\def\philscip{\philscisi}%\citeauthor{0philscicite}\addtocategory{ifarchcit}{0philscicite}%eprint
\texttt{PhilSci:\urlalt{http://philsci-archive.pitt.edu/archive/#1}{#1}}%
%\texttt{\href{http://philsci-archive.pitt.edu/archive/#1}{\protect\url{PhilSci:#1}}}%
%\providecommand{\mparcnote}{\texttt{philsci} eprints available at \url{http://philsci-archive.pitt.edu/}.}
}
\newcommand*{\biorxiveprint}[1]{%\global\def\biorxivp{\biorxivsi}%\citeauthor{0arxivcite}\addtocategory{ifarchcit}{0arxivcite}%eprint
\texttt{bioRxiv doi:\urlalt{https://doi.org/10.1101/#1}{10.1101/#1}}%
%\texttt{\href{http://arxiv.org/abs/#1}{\protect\url{arXiv:#1}}}%
%\renewcommand{\arxivnote}{\texttt{arXiv} eprints available at \url{http://arxiv.org/}.}
}
\newcommand*{\osfeprint}[1]{%
\texttt{Open Science Framework doi:\urlalt{https://doi.org/10.17605/osf.io/#1}{10.17605/osf.io/#1}}%
}

\usepackage{graphicx}

%\usepackage{wrapfig}

%\usepackage{tikz-cd}

\PassOptionsToPackage{hyphens}{url}\usepackage[hypertexnames=false]{hyperref}

\usepackage[depth=4]{bookmark}
\hypersetup{colorlinks=true,bookmarksnumbered,pdfborder={0 0 0.25},citebordercolor={0.2667 0.4667 0.6667},citecolor=mypurpleblue,linkbordercolor={0.6667 0.2 0.4667},linkcolor=myredpurple,urlbordercolor={0.1333 0.5333 0.2},urlcolor=mygreen,breaklinks=true,pdftitle={\pdftitle},pdfauthor={\pdfauthor}}
% \usepackage[vertfit=local]{breakurl}% only for arXiv
\providecommand*{\urlalt}{\href}

\usepackage[british]{datetime2}
\DTMnewdatestyle{mydate}%
{% definitions
\renewcommand*{\DTMdisplaydate}[4]{%
\number##3\ \DTMenglishmonthname{##2} ##1}%
\renewcommand*{\DTMDisplaydate}{\DTMdisplaydate}%
}
\DTMsetdatestyle{mydate}

%%%%%%%%%%%%%%%%%%%%%%%%%%%%%%%%%%%%%%%%%%%%%%%%%%%%%%%%%%%%%%%%%%%%%%%%%%%%
%%% Layout. I do not know on which kind of paper the reader will print the
%%% paper on (A4? letter? one-sided? double-sided?). So I choose A5, which
%%% provides a good layout for reading on screen and save paper if printed
%%% two pages per sheet. Average length line is 66 characters and page
%%% numbers are centred.
%%%%%%%%%%%%%%%%%%%%%%%%%%%%%%%%%%%%%%%%%%%%%%%%%%%%%%%%%%%%%%%%%%%%%%%%%%%%
\ifafour\setstocksize{297mm}{210mm}%{*}% A4
\else\setstocksize{210mm}{5.5in}%{*}% 210x139.7
\fi
\settrimmedsize{\stockheight}{\stockwidth}{*}
\setlxvchars[\normalfont] %313.3632pt for a 66-characters line
\setxlvchars[\normalfont]
\setlength{\trimtop}{0pt}
\setlength{\trimedge}{\stockwidth}
\addtolength{\trimedge}{-\paperwidth}
% The length of the normalsize alphabet is 133.05988pt - 10 pt = 26.1408pc
% The length of the normalsize alphabet is 159.6719pt - 12pt = 30.3586pc
% Bringhurst gives 32pc as boundary optimal with 69 ch per line
% The length of the normalsize alphabet is 191.60612pt - 14pt = 35.8634pc
\ifafour\settypeblocksize{*}{32pc}{1.618} % A4
%\setulmargins{*}{*}{1.667}%gives 5/3 margins % 2 or 1.667
\else\settypeblocksize{*}{26pc}{1.618}% nearer to a 66-line newpx and preserves GR
\fi
\setulmargins{*}{*}{1}%gives equal margins
\setlrmargins{*}{*}{*}
\setheadfoot{\onelineskip}{2.5\onelineskip}
\setheaderspaces{*}{2\onelineskip}{*}
\setmarginnotes{2ex}{10mm}{0pt}
\checkandfixthelayout[nearest]
\fixpdflayout
%%% End layout
%% this fixes missing white spaces
%\pdfmapline{+dummy-space <dummy-space.pfb}
%%\pdfglyphtounicode{space}{0020}
%\pdfinterwordspaceon%

%%% Sectioning
\newcommand*{\asudedication}[1]{%
{\par\centering\textit{#1}\par}}
\newenvironment{acknowledgements}{\section*{Thanks}\addcontentsline{toc}{section}{Thanks}}{\par}
\makeatletter\renewcommand{\appendix}{\par
  \bigskip{\centering
   \interlinepenalty \@M
   \normalfont
   \printchaptertitle{\sffamily\appendixpagename}\par}
  \setcounter{section}{0}%
  \gdef\@chapapp{\appendixname}%
  \gdef\thesection{\@Alph\c@section}%
  \anappendixtrue}\makeatother
\counterwithout{section}{chapter}
\setsecnumformat{\upshape\csname the#1\endcsname\quad}
\setsecheadstyle{\large\bfseries\sffamily%
\centering}
\setsubsecheadstyle{\bfseries\sffamily%
\raggedright}
%\setbeforesecskip{-1.5ex plus 1ex minus .2ex}% plus 1ex minus .2ex}
%\setaftersecskip{1.3ex plus .2ex }% plus 1ex minus .2ex}
%\setsubsubsecheadstyle{\bfseries\sffamily\slshape\raggedright}
%\setbeforesubsecskip{1.25ex plus 1ex minus .2ex }% plus 1ex minus .2ex}
%\setaftersubsecskip{-1em}%{-0.5ex plus .2ex}% plus 1ex minus .2ex}
\setsubsecindent{0pt}%0ex plus 1ex minus .2ex}
\setparaheadstyle{\bfseries\sffamily%
\raggedright}
\setcounter{secnumdepth}{2}
\setlength{\headwidth}{\textwidth}
\newcommand{\addchap}[1]{\chapter*[#1]{#1}\addcontentsline{toc}{chapter}{#1}}
\newcommand{\addsec}[1]{\section*{#1}\addcontentsline{toc}{section}{#1}}
\newcommand{\addsubsec}[1]{\subsection*{#1}\addcontentsline{toc}{subsection}{#1}}
\newcommand{\addpara}[1]{\paragraph*{#1.}\addcontentsline{toc}{subsubsection}{#1}}
\newcommand{\addparap}[1]{\paragraph*{#1}\addcontentsline{toc}{subsubsection}{#1}}

%%% Headers, footers, pagestyle
\copypagestyle{manaart}{plain}
\makeheadrule{manaart}{\headwidth}{0.5\normalrulethickness}
\makeoddhead{manaart}{%
{\footnotesize%\sffamily%
\scshape\headauthor}}{}{{\footnotesize\sffamily%
\headtitle}}
\makeoddfoot{manaart}{}{\thepage}{}
\newcommand*\autanet{\includegraphics[height=\heightof{M}]{autanet.pdf}}
\definecolor{mygray}{gray}{0.333}
\iftypodisclaim%
\ifafour\newcommand\addprintnote{\begin{picture}(0,0)%
\put(245,149){\makebox(0,0){\rotatebox{90}{\tiny\color{mygray}\textsf{This
            document is designed for screen reading and
            two-up printing on A4 or Letter paper}}}}%
\end{picture}}% A4
\else\newcommand\addprintnote{\begin{picture}(0,0)%
\put(176,112){\makebox(0,0){\rotatebox{90}{\tiny\color{mygray}\textsf{This
            document is designed for screen reading and
            two-up printing on A4 or Letter paper}}}}%
\end{picture}}\fi%afourtrue
\makeoddfoot{plain}{}{\makebox[0pt]{\thepage}\addprintnote}{}
\else
\makeoddfoot{plain}{}{\makebox[0pt]{\thepage}}{}
\fi%typodisclaimtrue
\makeoddhead{plain}{}{}{\footnotesize\reporthead}
% \copypagestyle{manainitial}{plain}
% \makeheadrule{manainitial}{\headwidth}{0.5\normalrulethickness}
% \makeoddhead{manainitial}{%
% \footnotesize\sffamily%
% \scshape\headauthor}{}{\footnotesize\sffamily%
% \headtitle}
% \makeoddfoot{manaart}{}{\thepage}{}

\pagestyle{manaart}

\setlength{\droptitle}{-3.9\onelineskip}
\pretitle{\begin{center}\Large\sffamily%
\bfseries}
\posttitle{\bigskip\end{center}}

\makeatletter\newcommand*{\atf}{\includegraphics[%trim=1pt 1pt 0pt 0pt,
totalheight=\heightof{@}]{atblack.png}}\makeatother
\providecommand{\affiliation}[1]{\textsl{\textsf{\footnotesize #1}}}
\providecommand{\epost}[1]{\texttt{\footnotesize\textless#1\textgreater}}
\providecommand{\email}[2]{\href{mailto:#1ZZ@#2 ((remove ZZ))}{#1\protect\atf#2}}

\preauthor{\vspace{-0.5\baselineskip}\begin{center}
\normalsize\sffamily%
\lineskip  0.5em}
\postauthor{\par\end{center}}
\predate{\DTMsetdatestyle{mydate}\vspace{0.5\baselineskip}%
\begin{center}\footnotesize}
\postdate{\end{center}\vspace{-\medskipamount}}

\setfloatadjustment{figure}{\footnotesize}
\captiondelim{\quad}
\captionnamefont{\footnotesize\sffamily%
}
\captiontitlefont{\footnotesize}
\firmlists*
\midsloppy
% handling orphan/widow lines, memman.pdf
% \clubpenalty=10000
% \widowpenalty=10000
% \raggedbottom
% Downes, memman.pdf
\clubpenalty=9996
\widowpenalty=9999
\brokenpenalty=4991
\predisplaypenalty=10000
\postdisplaypenalty=1549
\displaywidowpenalty=1602
\raggedbottom

\paragraphfootnotes\setlength{\footmarkwidth}{1em}
%\twocolumnfootnotes
%\setlength{\footmarksep}{0em}
\footmarkstyle{\textsuperscript{%\color{myred}
\scriptsize\bfseries#1}~}
%\footmarkstyle{\textsuperscript{[#1]}~}

\selectlanguage{british}\frenchspacing

%%%%%%%%%%%%%%%%%%%%%%%%%%%%%%%%%%%%%%%%%%%%%%%%%%%%%%%%%%%%%%%%%%%%%%%%%%%%
%%% Paper's details
%%%%%%%%%%%%%%%%%%%%%%%%%%%%%%%%%%%%%%%%%%%%%%%%%%%%%%%%%%%%%%%%%%%%%%%%%%%%
\title{\propertitle}
\author{%
\hspace*{\stretch{1}}%
\parbox{0.5\linewidth}%\makebox[0pt][c]%
{\protect\centering P.G.L.  Porta Mana\\%
\footnotesize\epost{\email{piero.mana}{ntnu.no}}}%
\hspace*{\stretch{1}}%
%% uncomment if additional authors present
\parbox{0.5\linewidth}%\makebox[0pt][c]%
{\protect\centering E. Torre\\%
\footnotesize\epost{\email{emiliano.torre}{yahoo.com}}}%
\hspace*{\stretch{1}}\\%
\parbox{0.5\linewidth}%\makebox[0pt][c]%
{\protect\centering V. Rostami\\%
\footnotesize\epost{\email{vrostami}{uni-koeln.de}}}%
\hspace*{\stretch{1}}%
\parbox{0.5\linewidth}%\makebox[0pt][c]%
{\protect\centering Y. Roudi\\%
\footnotesize\epost{\email{yasser.roudi}{ntnu.no}}}%
\hspace*{\stretch{1}}%
%\quad\href{https://orcid.org/0000-0002-6070-0784}{\protect\includegraphics[scale=0.16]{orcid_32x32.png}\textsc{orcid}:0000-0002-6070-0784}%
}

\date{Draft of \today\ (first drafted \firstdraft)}
%\date{\firstpublished; updated \updated}

%%%%%%%%%%%%%%%%%%%%%%%%%%%%%%%%%%%%%%%%%%%%%%%%%%%%%%%%%%%%%%%%%%%%%%%%%%%%
%%% Macros @@@
%%%%%%%%%%%%%%%%%%%%%%%%%%%%%%%%%%%%%%%%%%%%%%%%%%%%%%%%%%%%%%%%%%%%%%%%%%%%

% Common ones - uncomment as needed
%\providecommand{\nequiv}{\not\equiv}
%\providecommand{\coloneqq}{\mathrel{\mathop:}=}
%\providecommand{\eqqcolon}{=\mathrel{\mathop:}}
%\providecommand{\varprod}{\prod}
\newcommand*{\de}{\partialup}%partial diff
\newcommand*{\pu}{\piup}%constant pi
\newcommand*{\delt}{\deltaup}%Kronecker, Dirac
%\newcommand*{\eps}{\varepsilonup}%Levi-Civita, Heaviside
%\newcommand*{\riem}{\zetaup}%Riemann zeta
%\providecommand{\degree}{\textdegree}% degree
%\newcommand*{\celsius}{\textcelsius}% degree Celsius
%\newcommand*{\micro}{\textmu}% degree Celsius
\newcommand*{\I}{\mathrm{i}}%imaginary unit
\newcommand*{\e}{\mathrm{e}}%Neper
\newcommand*{\di}{\mathrm{d}}%differential
%\newcommand*{\Di}{\mathrm{D}}%capital differential
%\newcommand*{\planckc}{\hslash}
%\newcommand*{\avogn}{N_{\textrm{A}}}
%\newcommand*{\NN}{\bm{\mathrm{N}}}
%\newcommand*{\ZZ}{\bm{\mathrm{Z}}}
%\newcommand*{\QQ}{\bm{\mathrm{Q}}}
\newcommand*{\RR}{\bm{\mathrm{R}}}
%\newcommand*{\CC}{\bm{\mathrm{C}}}
%\newcommand*{\nabl}{\bm{\nabla}}%nabla
%\DeclareMathOperator{\lb}{lb}%base 2 log
%\DeclareMathOperator{\tr}{tr}%trace
%\DeclareMathOperator{\card}{card}%cardinality
%\DeclareMathOperator{\im}{Im}%im part
%\DeclareMathOperator{\re}{Re}%re part
%\DeclareMathOperator{\sgn}{sgn}%signum
%\DeclareMathOperator{\ent}{ent}%integer less or equal to
%\DeclareMathOperator{\Ord}{O}%same order as
%\DeclareMathOperator{\ord}{o}%lower order than
%\newcommand*{\incr}{\triangle}%finite increment
\newcommand*{\defd}{\coloneqq}
\newcommand*{\defs}{\eqqcolon}
%\newcommand*{\Land}{\bigwedge}
%\newcommand*{\Lor}{\bigvee}
%\newcommand*{\lland}{\DOTSB\;\land\;}
%\newcommand*{\llor}{\DOTSB\;\lor\;}
%\newcommand*{\limplies}{\mathbin{\Rightarrow}}%implies
%\newcommand*{\suchthat}{\mid}%{\mathpunct{|}}%such that (eg in sets)
%\newcommand*{\with}{\colon}%with (list of indices)
%\newcommand*{\mul}{\times}%multiplication
%\newcommand*{\inn}{\cdot}%inner product
%\newcommand*{\dotv}{\mathord{\,\cdot\,}}%variable place
%\newcommand*{\comp}{\circ}%composition of functions
%\newcommand*{\con}{\mathbin{:}}%scal prod of tensors
%\newcommand*{\equi}{\sim}%equivalent to 
\renewcommand*{\asymp}{\simeq}%equivalent to 
%\newcommand*{\corr}{\mathrel{\hat{=}}}%corresponds to
%\providecommand{\varparallel}{\ensuremath{\mathbin{/\mkern-7mu/}}}%parallel (tentative symbol)
\renewcommand*{\le}{\leqslant}%less or equal
\renewcommand*{\ge}{\geqslant}%greater or equal
\DeclarePairedDelimiter\clcl{[}{]}
%\DeclarePairedDelimiter\clop{[}{[}
%\DeclarePairedDelimiter\opcl{]}{]}
%\DeclarePairedDelimiter\opop{]}{[}
\DeclarePairedDelimiter\abs{\lvert}{\rvert}
%\DeclarePairedDelimiter\norm{\lVert}{\rVert}
\DeclarePairedDelimiter\set{\{}{\}}
%\DeclareMathOperator{\pr}{P}%probability
\newcommand*{\pf}{\mathrm{p}}%probability
\newcommand*{\p}{\mathrm{P}}%probability
\newcommand*{\E}{\mathrm{E}}
\renewcommand*{\|}{\nonscript\,\vert\nonscript\;\mathopen{}}
%\DeclarePairedDelimiterX{\cond}[2]{(}{)}{#1\nonscript\,\delimsize\vert\nonscript\;\mathopen{}#2}
\DeclarePairedDelimiterX{\condt}[2]{[}{]}{#1\nonscript\,\delimsize\vert\nonscript\;\mathopen{}#2}
%\DeclarePairedDelimiterX{\conds}[2]{\{}{\}}{#1\nonscript\,\delimsize\vert\nonscript\;\mathopen{}#2}
%\newcommand*{\+}{\lor}
%\renewcommand{\*}{\land}
\newcommand*{\sect}{\S}% Sect.~
\newcommand*{\sects}{\S\S}% Sect.~
\newcommand*{\chap}{ch.}%
\newcommand*{\chaps}{chs}%
\newcommand*{\bref}{ref.}%
\newcommand*{\brefs}{refs}%
%\newcommand*{\fn}{fn}%
\newcommand*{\eqn}{eq.}%
\newcommand*{\eqns}{eqs}%
\newcommand*{\fig}{fig.}%
\newcommand*{\figs}{figs}%
\newcommand*{\vs}{{vs}}
%\newcommand*{\etc}{{etc.}}
%\newcommand*{\ie}{{i.e.}}
%\newcommand*{\ca}{{c.}}
\newcommand*{\eg}{{e.g.}}
\newcommand*{\foll}{{ff.}}
%\newcommand*{\viz}{{viz}}
\newcommand*{\cf}{{cf.}}
%\newcommand*{\Cf}{{Cf.}}
%\newcommand*{\vd}{{v.}}
\newcommand*{\etal}{{et al.}}
%\newcommand*{\etsim}{{et sim.}}
%\newcommand*{\ibid}{{ibid.}}
%\newcommand*{\sic}{{sic}}
%\newcommand*{\id}{\mathte{I}}%id matrix
%\newcommand*{\nbd}{\nobreakdash}%
%\newcommand*{\bd}{\hspace{0pt}}%
%\def\hy{-\penalty0\hskip0pt\relax}
\newcommand*{\labelbis}[1]{\tag*{(\ref{#1})$_\text{r}$}}
%\newcommand*{\mathbox}[2][.8]{\parbox[t]{#1\columnwidth}{#2}}
%\newcommand*{\zerob}[1]{\makebox[0pt][l]{#1}}
\newcommand*{\tprod}{\mathop{\textstyle\prod}\nolimits}
\newcommand*{\tsum}{\mathop{\textstyle\sum}\nolimits}
%\newcommand*{\tint}{\begingroup\textstyle\int\endgroup\nolimits}
%\newcommand*{\tland}{\mathop{\textstyle\bigwedge}\nolimits}
%\newcommand*{\tlor}{\mathop{\textstyle\bigvee}\nolimits}
%\newcommand*{\sprod}{\mathop{\textstyle\prod}}
%\newcommand*{\ssum}{\mathop{\textstyle\sum}}
%\newcommand*{\sint}{\begingroup\textstyle\int\endgroup}
%\newcommand*{\sland}{\mathop{\textstyle\bigwedge}}
%\newcommand*{\slor}{\mathop{\textstyle\bigvee}}
\newcommand*{\T}{^\intercal}%transpose
%%\newcommand*{\QEM}%{\textnormal{$\Box$}}%{\ding{167}}
%\newcommand*{\qem}{\leavevmode\unskip\penalty9999 \hbox{}\nobreak\hfill
%\quad\hbox{\QEM}}

%%%%%%%%%%%%%%%%%%%%%%%%%%%%%%%%%%%%%%%%%%%%%%%%%%%%%%%%%%%%%%%%%%%%%%%%%%%%
%%% Custom macros for this file @@@
%%%%%%%%%%%%%%%%%%%%%%%%%%%%%%%%%%%%%%%%%%%%%%%%%%%%%%%%%%%%%%%%%%%%%%%%%%%%
\definecolor{notecolour}{RGB}{68,170,153}
\newcommand*{\puzzle}{{\fontencoding{U}\fontfamily{fontawesometwo}\selectfont\symbol{225}}}
%\newcommand*{\puzzle}{\maltese}
\newcommand{\mynote}[1]{ {\color{notecolour}\puzzle\ #1}}
\newcommand*{\widebar}[1]{{\mkern1.5mu\skew{2}\overline{\mkern-1.5mu#1\mkern-1.5mu}\mkern 1.5mu}}

% \newcommand{\explanation}[4][t]{%\setlength{\tabcolsep}{-1ex}
% %\smash{
% \begin{tabular}[#1]{c}#2\\[0.5\jot]\rule{1pt}{#3}\\#4\end{tabular}}%}
% \newcommand*{\ptext}[1]{\text{\small #1}}
%\DeclareMathOperator*{\argsup}{arg\,sup}
\newcommand*{\dob}{degree of belief}
\newcommand*{\dobs}{degrees of belief}

\newcommand*{\tav}{\widehat} %time average
\newcommand*{\av}{\widebar} %pop average
\newcommand*{\sav}{\widebar} %subpop average

\newcommand*{\yG}{G}

%\newcommand*{\yXv}{\varSigma}
\newcommand*{\yAv}{A}
%\newcommand*{\yxv}{s}
\newcommand*{\yav}{a}
\newcommand*{\yNg}{N_{\text{g}}}
\newcommand*{\yNng}{N_{\text{ng}}}
\newcommand*{\yNN}{\varNu}
\newcommand*{\yFF}{F}
\newcommand*{\yff}{f}
\newcommand*{\yF}{\bm{\yFF}}
\newcommand*{\yf}{\bm{\yff}}
\newcommand*{\yFFm}{\yFF^{*}}
\newcommand*{\yFm}{\yF^{*}}

%\newcommand*{\yx}{\bm{\yxv}}%subpop state
%\newcommand*{\yxs}{\sav{\yx}}%subpop av state
%\newcommand*{\yX}{\bm{\yXv}}%pop state
%\newcommand*{\yXf}{\av{\yX}}%pop av state
\newcommand*{\yaa}{\bm{\yav}}%subpop value
\newcommand*{\yaas}{\sav{\yaa\yaa}}%subpop av state
%\newcommand*{\yXXf}{\av{\yAA\yAA}}%subpop av state
\newcommand*{\ya}{\yav}%subpop av value
\newcommand*{\yAA}{\bm{\yAv}}%pop value
\newcommand*{\yA}{\yAv}%pop av value
%conditional assumptions
\newcommand*{\yH}{\varIota}
\newcommand*{\yHa}{\varIota_\textrm{p}}
\newcommand*{\yHb}{\varIota_\textrm{s}}
\newcommand*{\yHc}{\varKappa}
\newcommand*{\yHd}{\varIota_\textrm{u}}
\newcommand*{\yHp}{\varIota'}
\newcommand*{\yHi}{\varIota''}
\newcommand*{\yI}{\varIota}
\newcommand*{\yD}{D}
%Lagr multipliers
\newcommand*{\yg}{\bm{c}}
\newcommand*{\yc}{\Hat{\bm{c}}}
\newcommand*{\yll}{\bm{\lambda}}
\newcommand*{\yl}{\lambda}
\newcommand*{\yk}{z}
\newcommand*{\yK}{\zeta}
%entropies
\newcommand*{\re}{\Eta}
\newcommand*{\ybu}{\re_{\text{B}}}

\newcommand*{\yr}{\bm{r}}

\newcommand*{\ynuu}{\nu}
\newcommand*{\ynu}{\bm{\ynuu}}
\newcommand*{\yjj}{\phi}
\newcommand*{\yj}{\bm{\yjj}}

\newcommand*{\prop}[1]{`#1'}

\newcommand{\appropto}{\mathrel{\vcenter{
  \offinterlineskip\halign{\hfil$##$\cr
    \propto\cr\noalign{\kern2pt}\sim\cr\noalign{\kern-2pt}}}}}
% https://tex.stackexchange.com/a/33547/97039
%\newcommand*{\fm}{f}
\newcommand*{\px}{P_{\text{me}}}
\newcommand*{\pxx}{p}
\newcommand*{\pxxx}{P}
\newcommand*{\Nd}{\textrm{N}}
\newcommand*{\ycmm}{C}
\newcommand*{\ycm}{\mathte{\ycmm}}
\newcommand*{\ycs}{\bm{\sigma}}
\newcommand*{\yL}{L}

\newcommand*{\yla}{\bm{\lambda}}
\newcommand*{\yJ}{\bm{J}}
%\newcommand*{\yd}{\bm{\varDelta}}
\newcommand*{\ydi}{\varDelta}

\newcommand*{\mee}{maximum entropy}
\newcommand*{\me}{maximum-entropy}

\newcommand*{\sh}{\mathrm{\varEta}}
\newcommand*{\nat}{\textrm{nat}}
\newcommand*{\bit}{\textrm{bit}}
\newcommand*{\hart}{\textrm{Hart}}

\newsubfloat{figure}
\subcaptionsize{\normalfont}
\subcaptionlabelfont{\normalfont}
\subcaptionfont{\normalfont}
%%% Custom macros end @@@

%%%%%%%%%%%%%%%%%%%%%%%%%%%%%%%%%%%%%%%%%%%%%%%%%%%%%%%%%%%%%%%%%%%%%%%%%%%%
%%% Beginning of document
%%%%%%%%%%%%%%%%%%%%%%%%%%%%%%%%%%%%%%%%%%%%%%%%%%%%%%%%%%%%%%%%%%%%%%%%%%%%
\firmlists
\raggedbottom
\begin{document}
\captiondelim{\quad}\captionnamefont{\footnotesize}\captiontitlefont{\footnotesize}
\selectlanguage{british}\frenchspacing
\maketitle

%%%%%%%%%%%%%%%%%%%%%%%%%%%%%%%%%%%%%%%%%%%%%%%%%%%%%%%%%%%%%%%%%%%%%%%%%%%%
%%% Abstract
%%%%%%%%%%%%%%%%%%%%%%%%%%%%%%%%%%%%%%%%%%%%%%%%%%%%%%%%%%%%%%%%%%%%%%%%%%%%
\abstractrunin
\abslabeldelim{}
\renewcommand*{\abstractname}{}
\setlength{\absleftindent}{0pt}
\setlength{\absrightindent}{0pt}
\setlength{\abstitleskip}{-\absparindent}
\begin{abstract}\labelsep 0pt%
\noindent This work shows how to build a maximum-entropy probabilistic
model for the total activity of a large population of unrecorded neurons,
given only some activity data or statistics recorded from a small sample
thereof. The model addresses only the collective behaviour of the larger
population, thus gaining in simplicity and computational effectiveness at
the expense of more detailed inferential power.

The model is applied to two sets of neuronal data from the literature. It
makes interesting forecasts about the larger population, such as the
probable presence of two low-regime modes in the total activity, which are
not present in the sample data and could not be found with maximum-entropy
models applied only to the sample.

Using the same data sets it is shown how to use the model for questions
about the \enquote{cooperativity}, \enquote{interaction}, or
\enquote{synchrony} of neuronal activity. These questions are translated
into a hypothesis-testing problem about informational sufficiency, and
effectively answered with Bayes's theorem. The answers given by the present
model can be at variance from those given by a maximum-entropy model at the
sample level only. The reason is discussed: informational sufficiency is not
preserved under sampling.

The question of how large the larger population addressed by model should
or can be, and the assumptions on which the model is based, from a Bayesian
point of view, are also discussed.
% \\\noindent\emph{\footnotesize Note:
%   Dear Reader \amp\ Peer, this manuscript is being peer-reviewed by you.
%   Thank you.}
% \par%\\[\jot]
% \noindent
% {\footnotesize PACS: ***}\qquad%
% {\footnotesize MSC: ***}%
%\qquad{\footnotesize Keywords: ***}
\end{abstract}
\selectlanguage{british}\frenchspacing

%%%%%%%%%%%%%%%%%%%%%%%%%%%%%%%%%%%%%%%%%%%%%%%%%%%%%%%%%%%%%%%%%%%%%%%%%%%%
%%% Epigraph
%%%%%%%%%%%%%%%%%%%%%%%%%%%%%%%%%%%%%%%%%%%%%%%%%%%%%%%%%%%%%%%%%%%%%%%%%%%%
% \asudedication{\small ***}
% \vspace{\bigskipamount}
% \setlength{\epigraphwidth}{.7\columnwidth}
% %\epigraphposition{flushright}
% \epigraphtextposition{flushright}
% %\epigraphsourceposition{flushright}
% \epigraphfontsize{\footnotesize}
% \setlength{\epigraphrule}{0pt}
% %\setlength{\beforeepigraphskip}{0pt}
% %\setlength{\afterepigraphskip}{0pt}
% \epigraph{\emph{text}}{source}



%%%%%%%%%%%%%%%%%%%%%%%%%%%%%%%%%%%%%%%%%%%%%%%%%%%%%%%%%%%%%%%%%%%%%%%%%%%%
%%% BEGINNING OF MAIN TEXT
%%%%%%%%%%%%%%%%%%%%%%%%%%%%%%%%%%%%%%%%%%%%%%%%%%%%%%%%%%%%%%%%%%%%%%%%%%%%

% \mynote{comment about the possibility of drawing conclusions about a brain
%   area using different sets of neurons (eg because of recording across many
%   sessions)}


\section{Introduction:\\a model for collective inferences about large neuronal
  populations}
\label{sec:intro}

What correlations are dominant in the neuronal activity of a specific brain
area? How does such activity change when external stimuli or the activity
of other areas change? Does such activity range over all its mathematically
possible values, or only within restricted bounds?

Answering this kind of questions always engages an element of uncertainty.
We cannot say \enquote{the answer is such and such}; at best we can assign
degrees of reasonable belief -- that is, probabilities -- to the possible
answers. The assessment of such distributions of belief involves
experimental data, such as recordings of neuronal activity from specific
brain areas, and pre-data knowledge and hypotheses about biological
conditions and mechanisms. When we translate such pre-data beliefs into an
initial probability distribution we often simplify them in more or less
realistic ways, to make the analysis mathematically tractable. This is why
such initial probability distributions are called \enquote{(probabilistic)
  models}; and that is how we intend this word in the present work.

The best experimental measurements of instantaneous neuronal activity use
remarkable technologies \citep{hongetal2019}, but can still only record a
very small sample of neurons -- hundreds at most -- compared to the numbers
that constitute a functionally distinct brain region. Many probabilistic
models focus on such samples only: they neglect that the recorded neurons
are a sample from a larger population. Some probabilistic models try to
take unrecorded neurons into account, but they do so by describing each
neuron individually, thus becoming very complex
\citep{pillowetal2008,dunnetal2013,tyrchaetal2014,huang2015,battistinetal2017}.
It would be useful to explore models that operate in between, addressing
the activity of the larger brain area whence the sample comes, but
\emph{collectively}, that is, without asking about individual neuronal
details. Such intermediate models would be useful at least for preliminary
investigations, to help us decide which hypotheses to discard, which to
consider for a deeper analysis with more complex and costly models, or to
suggest new hypotheses.

% \mynote{[Yasser:] In the present work we aim to address the issue of
%   building probability distributions over a large population using
%   recording from subpopulations by considering the distribution of total
%   population activity as an example. [} 
In the present work we propose such an intermediate probabilistic model
\citep[\cf][\sect~3]{portamanaetal2015}. It answers this question: how much was
the \emph{total} activity of a large neuronal population, given the
observation of the activity of a very small sample thereof? This model
addresses a larger brain area, avoiding the assumption of isolation of the
sample; and by focusing on the total activity, rather than the activity of
individual neurons, it remains simple and numerically tractable.


The model can be viewed as a straightforward combination of the
maximum-entropy method and basic sampling relations from the
probability-calculus, discussed in \sect~\ref{sec:model}. The
maximum-entropy or minimum-relative-entropy method
\citep{jaynes1957}[\chap~11]{jaynes1994_r2003}[for the minimum-relative
entropy method see:][\sect~4.b]{jaynes1963}[for a more rigorous
derivation]{hobsonetal1973}[for applications]{sivia1996_r2006}
%[see also][for applications to irreversible statistical mechanics]{grandy1980}
to build initial probability distributions has been used in the
neurosciences for several kinds of investigations about the neuronal
activity of various brain areas and about other phenomena of importance % ,
% for example gene and protein interaction
\citep[]{martignonetal1995,bohteetal2000,shlensetal2006,schneidmanetal2006,tkaciketal2006,roudietal2009c,barreiroetal2010,shimazakietal2015,moraetal2015%,lezonetal2006,weigtetal2009
}[for
further references see][]{lathametal2013}. However, our model can also be
viewed as an approximation of a more detailed model within Bayesian theory;
we shall constantly keep this point of view in mind in our discussion.

We illustrate possible uses of our model in \sect~\ref{sec:application}, by
applying it to two concrete data sets:
\begin{enumerate*}[label=(\alph*)]
\item the activity of 65 neurons recorded for 20 min from a rat's Medial
  Entorhinal Cortex \citep{stensolaetal2012},
\item the activity of 159 neurons recorded for 15 min from a macaque's
  Motor Cortex \citep{rostamietal2016_r2017}.
\end{enumerate*}
For each data set the model gives us the most plausible frequencies of all
levels of total activity of a much larger population, 1\,000--10\,000
neurons in size, during the recording. For example it can tell us that 60
out of 10\,000 neurons were most likely active during 1\% of the recording
time (though not necessarily always the same 60), 250 neurons out of
10\,000 were active during 0.4\% of the recording time, and so on. The
precise meaning of this frequency distribution is explained in
\sect~\ref{sec:model}. For the two example data sets, the guessed frequency
distributions for the larger population are distinctly different from those
for the sample. For example, for the first data set the frequency
distribution for the larger population has two very distinct modes, both at
low activities %($0.7\%$ and $2.6\%$ of simultaneously active neurons
% within $3\;\mathrm{ms}$,
(see \fig~\ref{fig:plots_5mom}), whereas the frequency distribution for the
sample is monotonically decreasing with its maximum at zero activity. These
results show that the proposed model can be used for the formulation or
preliminary assessment of interesting hypotheses. We illustrate this use
with toy examples in \sects~\ref{sec:application}
and~\ref{sec:marginalization}. Note that it is not possible to guess these
features of the larger population by applying the maximum-entropy method
\emph{at the level of the sample alone}.

% We want to stress the usefulness of making a quantified guess about the
% activity of a larger brain area. Such a guess seems indeed to be the
% primary idea behind recording a sample from that area. %  Second, it allows us
% % to make comparisons across experimental sessions; such comparisons would me
% % difficult or meaningless if made with the recorded samples, which generally
% % comprise non-overlapping sets of neurons and differ in size.
% There's also an important methodological reason.

Our model also solves a methodological problem in the use of
maximum-entropy methods to assess the \enquote{cooperativity},
\enquote{interaction}, or \enquote{synchrony} in neuronal activity, for
example studying its pairwise correlations and correlations of higher order
\citep[see for
example][]{martignonetal1995,bohteetal2000,schneidmanetal2006,shlensetal2006,barreiroetal2010,ganmoretal2011,granotatedgietal2013}.
When the difference in size between a large population and a sample thereof
is too large, the sufficiency of correlations of some order for the sample
implies the \emph{lack} of sufficiency of correlations of the same order
for the larger population, and vice versa. Maximum-entropy applications at
the sample level can therefore deceive us in questions regarding the
cooperativity of the larger population. Our model addresses directly the
latter instead. We explain this point more precisely in
\sect~\ref{sec:marginalization}.


% Maximum-entropy models have often been used to test whether a smaller part
% of the sample data, for example time averages and second-order
% correlations, is sufficient for quantifying our \dob\ about some aspects of
% the neuronal activity of the sample. This use of maximum-entropy at the
% sample level, however, has an important issue. We discuss this issue and
% explain how our approach bypasses it in \sect~\ref{sec:marginalization}.
% Moreover, our proposed model gives clearer answers to questions of
% correlation sufficiency: we show this by applying it to the two data set
% above addressing an example question of sufficiency, two moments versus
% four moments.

How large is the \enquote{larger population} considered by the model
proposed here? It obviously cannot include the full brain. The size of the
larger population is determined by the validity of the formulae from
sampling theory, and can range from a set of neurons around the recording
probe to part of a brain area, depending on some assumptions about the
recording procedure and the brain region under consideration. This matter
is discussed in detail in \sect~\ref{sec:N}.

In \sect~\ref{sec:assumptions} we discuss in detail the assumptions and
approximations which define our model, from the point of view of the
probability-calculus.

% \mynote{review this}We obviously don't know whether the activity levels of
% the larger population really had the frequencies given by the model, during
% the recording. The model gives our most plausible guess. In high
% dimensions, however, the features of the most plausible distribution may
% \emph{not} be typical of the majority of most plausible distributions; and
% the set of all possible frequency distributions, if the full population for
% example comprises 1\,000 neurons, is a 1\,000-dimensional space. In
% \sect~\ref{sec:assumptions} we therefore try to assess which features of
% the frequency distribution delivered by our method may be typical and
% therefore expected of the actual one. We find that general features such as
% the bimodality of the first data set are indeed typical. Maximum-entropy
% models can be considered as approximations of Bayesian models based on
% various assumptions of inferential sufficiency
% \citep{jaynes1986d_r1996,portamana2017}. How do our guesses change if we
% modify our pre-data assumptions? We show, in the same section, that the
% typical features indicated by our method are robust against such changes.

A summary of all points above and a discussion of the usefulness of the
model is given in the final \sect~\ref{sec:summary_discussion}.

\medskip


Our notation and terminology follow \textsc{iso}
\citey{iso2006,iso2006b,jcgm1993_r2008} standards and Jaynes
\citey{jaynes1994_r2003} for probability. We use \enquote{\dob},
\enquote{belief}, and \enquote{probability} interchangeably.




% (for a full population of 1\,000 neurons and
% a recording of 500\,000 time bins, for example, there are
% $\binom{500\,000+1\,000}{1\,000} \approx 10^{3\,131}$ possible frequency
% distributions).

% Later on we'll show that our \dobs\ can be quantified by
% exclusively using the probability-calculus. This derivation will provide a
% more accurate quantification, revealing that the maximum-entropy answer is
% only a first approximation.


% \mynote{Say something more about advantages of such question/answer: for
%   example we can make statements about total activity of brain area even
%   across recordings when sampled neurons aren't the same.}

% \mynote{Add this?: from sampling theory we know that important
%   features of the full population may not be visible in a sample because
%   smoothed out. But using sampling theory in the inverse direction we can
%   infer such full-population features from the sample.}

% \mynote{To be continued after structure of the rest of the article is
%   clear. Orig intro is on p.~\pageref{orig_intro}}



\section{Model: maximum-entropy and sampling,\\ and a Bayesian perspective}
\label{sec:model}

Some context and mathematical notation first.

The context we consider is as follows. During an experimental session we
have recorded the spiking activities of $n$ neurons for a certain amount of
time. These neurons are our \enquote{sample}. Their spikes are binned into
$T$ time bins, and a binary value 1 is assigned to each bin containing at
least one spike, 0 to the bins without spikes. Call $\ya_{t}$ the number of
neurons that fire during time bin $t$: this is the \emph{total activity} of
the sample, or just \enquote{activity} for short. Obviously
$\ya_{t} \in\set{0,1,\dotsc,n}$; if $\ya_{t}=0$, no neuron spikes during
bin $t$; if $\ya_{t}=n$, all spike (at least once) at some point during bin
$t$, and so on. For brevity we say \enquote{at $t$} for \enquote{during
  time bin
  $t$}. % If we divide the total activity by the population size we have the
% \emph{normalized total activity} or population-averaged activity $\ya/n$,
% ranging from $0$ to $1$ in $1/n$ steps.
From the sequence of activities
$(\ya_{t})\equiv (\ya_{1}, \dotso, \ya_{T})$ we can count in how many bins
each activity level $\ya=0$, $\ya=1$, \ldots, $\ya=n$ appeared during the
recording. Dividing each number of bins by the total number of bins $T$ we
obtain the distribution of measured relative frequencies
$(\yff_{0},\dotso, \yff_{n}) \equiv (\yff_{\ya}) \defs \yf$. We can also
consider the (unknown) sample activity at time bins \emph{outside} of the
recorded period.


For many animal species, the neurons that are recorded within a brain area
are not specifically chosen from among the rest, owing to several limiting
factors; for example, limitations in how precisely electrodes are inserted.
The sample of recorded neurons may even change slightly across experimental
sessions that are very far apart in time. We assume that there's a volume,
comprising a population of $N$ neurons, for which we believe that any other
sample of size $n$ could have equally plausibly been recorded instead of
the sample of $n$ neurons that was actually recorded. This is what we will
mean with \enquote{larger population}. This population need not be a whole
functionally or anatomically distinct region. Loosely speaking it is the
area of which we believe our sample to be \enquote{representative}
\citep[we intend this term in the Meaning~2 of][]{kruskaletal1979b}[or
Meanings~8 and~9 of][]{kruskaletal1979c}[note that the
standard][\sect~3.1.14 states \enquote{The notion of representative sample
  is fraught with controversy, with some survey practitioners rejecting the
  term altogether}]{iso2014}.

The total activity of the larger population of $N$ neurons at $t$ is
$\yA_{t}$. The relative frequencies of the various activity levels during
the recording were
$(\yFF_{0},\dotso,\yFF_{N}) \equiv (\yFF_{\yA}) \defs \yF$. The values
$\yA_{t}$ at each $t$ and the frequency distribution $\yF$ are unknown to
us. We only know for certain that $\yA_{t} \in \set{0, 1, \dotsc, N}$, that
$\yA_{t} \ge \ya_{t}$, and that $N-\yA_{t} \ge n-\ya_{t}$ for obvious
reasons. For the time being we assume that we know $N$; in
\sect~\ref{sec:N} we discuss the consequences of our lack of precise
knowledge about this number.

Our questions concern general features of the total activity $\yA$ of the
larger population during and after the recording, and across sessions under
the same study conditions. For example: what was its frequency distribution
$\yF$ during the recording? How much does this frequency distribution
change across sessions? How much total activity should we expect at any
time bin during a recording? The model presented here gives a probability
distribution over the answers to these questions.


\bigskip

The idea behind our approach can be easily summarized in terms of the
maximum-entropy method:
\begin{enumerate}[label=(\alph*)]
\item\label{item:sample_step} Using sampling theory we determine the
  relation between some expected values -- specifically, moments -- for the
  total activity $\ya$ of the sample and corresponding expected values %$M$
  for the total activity $\yA$ of the larger population.
\item\label{item:maxent_step} Using the maximum-entropy method we build a
  distribution $\px(\yA \| M, N)$ for the total activity of the larger
  population of $N$ neurons, using as constraints the expected values $M$
  found in the previous step.
\end{enumerate}
We now discuss these steps more in detail, but leave their precise
mathematical implementation and a more detailed list of references to
appendix~\ref{sec:derivation_maxent_fullnet}.

\medskip

Step~\ref{item:sample_step} is just an application of the
probability-calculus, which gives an exact linear relation between the
first $m$ moments for the larger population and the first $m$ for the
sample \citep[see][for a complete discussion and proofs. Some proofs are
summarized again here in
appendix~\ref{sec:derivation_maxent_fullnet}]{portamanaetal2015}. The ones
determine the others and vice versa at every time bin. This relation holds
for any belief distribution $\pxxx(\yA_{t})$ for the larger-population
activity and its marginal $\pxx(\ya_{t})$ for the sample activity at that
bin. These two distributions are related by
\begin{equation}
  \label{eq:conditional_hypergeometric_t}
  \pxx(\ya_{t})= \sum_{\yA_{t}=0}^{N} \yG_{\ya_{t}\,\yA_{t}} \;\pxxx(\yA_{t})
  \qquad \ya_{t} \in \set{0,\dotsc,n},
\end{equation}
where $\yG_{\ya\yA}$ is the hypergeometric distribution defined in
\eqn~\eqref{eq:conditional_hypergeometric}, characteristic of
\enquote{drawing without replacement}
\citep[\chap~3]{jaynes1994_r2003}.%[\sect~4.8.3]{ross1976_r2010}[\sect~II.6]{feller1950_r1968}.


This sampling relation is even more straightforward if instead of power
moments we use \emph{normalized factorial moments} \citep{potts1953}. The
$m$th normalized factorial moment of a distribution $\pxx(\ya)$ is
\begin{equation}
  \label{eq:def_factorial_moment_a}
  % \E\Biggl[\binom{\ya}{m} \Biggr]\bigg/\binom{n}{m}
  % \defd
  \sum_{\ya=0}^{n} %m!\,
  \displaystyle\binom{\ya}{m} \bigg/ \binom{n}{m}\; \pxx(\ya)
  \qquad m \in \set{1,\dotsc,n},
\end{equation}
% where $\px(\ya \| \yH)$ is the probability for activity $\ya$ given
% assumptions $\yH$; \enquote{$\ya$} may refer to any recorded bin or to a
% new bin.
that is, the expected number of distinct $m$-tuples of simultaneously
active neurons (within a time-bin's width), normalized by the maximum
possible number of distinct $m$-tuples. % For
% example, with $m=2$, if $a=4$ neurons spike in a population of $n=5$, we
% have $\tbinom{4}{2}=6$ distinct pairs of simultaneously spiking neurons,
% and the total number of distinct pairs is $\tbinom{5}{2}=10$. The
% normalized number of spiking pairs is therefore $6/10$.
Note that the first $m$ factorial moments together provide the same
information as the first $m$ power moments together, and vice versa: they
are linearly related because $\tbinom{\ya}{m}$ is a polynomial in $\ya$ of
degree $m$. So we'll just say \enquote{first $m$ moments} from now
on. % Specifying one set is therefore
% equivalent to specifying the other set.
But the normalized factorial moments have a mathematically convenient
property: \emph{the first $n$ normalized factorial moments for the sample
  and for the larger population are numerically identical}:\citep[Appendix]{portamanaetal2015}
\begin{equation}
  \label{eq:equal_m_factorialmoments}
  % \begin{split}
  % \E\Biggl[\binom{\ya}{m} \Biggr]\bigg/\binom{n}{m}
  % &=
  % \E\Biggl[\binom{\yA}{m} \Biggr]\bigg/\binom{N}{m}
  % \qquad\text{or}\\
    \sum_{\ya=0}^{n} %m!\,
  \binom{\ya}{m}\bigg/  \binom{n}{m}\; \pxx(\ya)
=    \sum_{\yA=0}^{N} %m!\,
  \binom{\yA}{m}\bigg/  \binom{N}{m}\; \pxxx(\yA),\qquad
  % \E\condt[\big]{\tbinom{\ya}{m}\;\tbinom{n}{m}^{-1}}{\yH} =
  % \E\condt[\big]{\tbinom{\yA}{m}\;\tbinom{N}{m}^{-1}}{\yH},
  m \in \set{1,\dotsc,n}.
%\end{split}
\end{equation}
% where $\pxxx(\yA)$ is our distribution of belief about the larger-population
% activity, and $\pxx$


\medskip


In step~\ref{item:maxent_step} we actually use the
\emph{minimum-relative-entropy} method
\citep{hobsonetal1973,csiszar1985}[\sect~5.2.2]{sivia1996_r2006} with
respect to a uniform reference distribution. We shall still call it
\enquote{maximum-entropy} for brevity. It amounts to two prescriptions:
first, take the distributions satisfying specific convex constraints --
such as fixed expectations -- and among them select the one having minimum
relative entropy with respect to a reference distribution; second, judge --
and therefore set -- those expectations to be equal to some measured
averages, typically time averages. (A judgement is required because
probability expectations and empirical averages differ numerically in
general; a misjudgement of their approximate equality can lead to
\enquote{silly answers}, to quote MacKay \citep[p.~308]{mackay1995_r2003}.)

In our case we don't know the time averages of the quantity
$\binom{\yA}{m}$, so we cannot directly equate them to the factorial moment
$\E\bigl[\binom{\yA}{m}\bigr]$ of the larger-population distribution
$\pxxx(\yA)$. But \eqn~\eqref{eq:equal_m_factorialmoments} of
step~\ref{item:sample_step} comes to our rescue, because it says that the
expectation for the larger population are determined by that for the sample
$\E\bigl[\binom{\ya}{m}\bigr]$, and we do have the time average of its
corresponding sample quantity $\binom{\ya}{m}$. So we can combine the two
steps:
\begin{equation*}
  \label{eq:chain_prescription_moments}
  \mathrlap{\overbracket{\phantom{\text{measured averages} \rightarrow
        \text{sample moments}}}^{\textit{maximum-entropy prescription}}}
  \text{measured averages} \rightarrow
        \underbracket{\text{sample moments}
  \rightarrow
    \text{larger-population moments}}_{\textit{sampling theory}}
\end{equation*}
We obtain a distribution $\px(\yA \| M, N)$ for the larger population of
$N$ neurons by constraining some of its factorial moments, for example
$M=\set{1,\dotsc,m'}$ with $m'\le n$, to be equal to the sample's recorded
averages. In formulae, the constraints on $\pxxx(\yA)$ are
\begin{equation}
  \label{eq:final_constraint}
  \underbracket{\frac{1}{T}\sum_{t}\binom{\ya_{t}}{m}\bigg/\binom{n}{m}
  \equiv
   \sum_{\ya}
  \binom{\ya}{m}\bigg/  \binom{n}{m} \; \yff_{\ya}}_{\text{measured averages}}
  \mathrel{\;\;\underset{\mathclap{\overset{\Bigg\uparrow}{\text{our judgement}}}}{=}\;\;}
   \underbracket{\sum_{\yA}
     \binom{\yA}{m}\bigg/  \binom{N}{m} \; \pxxx(\yA)}_{\text{distribution moments}}
   \qquad m\in M.
\end{equation}
\smallskip

The result is the distribution of the form
\begin{equation}
  \label{eq:app_maxent_pop_first}
  \begin{gathered}
  \px(\yA \| M, N)  = \frac{1}{Z(\yll)}\;
%  g(\yA)\;
  \exp\Biggl[\sum_{m}\yl_{m}\;
  \binom{\yA}{m}\bigg/\binom{N}{m}
  \Biggr]
  \\
\text{with}\qquad  Z(\yll)\defd \sum_{\yA}
  % g(\yA)\;
  \exp\Biggl[\sum_{m}\yl_{m}\;
  \binom{\yA}{m}\bigg/\binom{N}{m}
  \Biggr],
\end{gathered}
\end{equation}
where the $m'$ parameters
$\yll \defd (\yl_{m}) \equiv(\yl_{1},\dotso,\yl_{m'})$ are determined by
the constraints~\eqref{eq:final_constraint} and can be found by convex
extremization, discussed further in
appendix~\ref{sec:derivation_maxent_fullnet}. Note that for $N>n$ the
constraint equation~\eqref{eq:final_constraint} may never be exactly
satisfied by the distribution~\eqref{eq:app_maxent_pop_first} for any
$\yla$. In this case the extremization minimizes the discrepancy between
the left and right sides of~\eqref{eq:final_constraint}. This discrepancy
comes from an approximation implicit in the maximum-entropy method (namely
that the number of time bins $T$ is infinite) combined with one assumption
specific to our application: that an activity level $\yA$ can equally
likely be generated by any set of $\yA$ neurons in the larger population.
The magnitude of this discrepancy can be a signature of the presence of
neuronal \enquote{assemblies} or \enquote{clusters}
\citep[\chap~12]{gerstneretal2014}{hebb1949_r2002}. We discuss this matter
in \sect~\ref{sec:assumptions}.

\bigskip

One important question remains about the distribution $\px$ of
\eqn~\eqref{eq:app_maxent_pop_first}: \emph{what} is it a distribution of?
-- Of probability? Of relative frequency? The answer is related to an
important relation between maximum-entropy distributions and the
probability-calculus, and also to our alternative point of view of the
present model.

% Before applying the formula above to two concrete data sets we want to add
% two remarks about the maximum-entropy method that are seldom made in the
% neuroscientific literature. They are important for the interpretation of
% the results and are further discussed in \sect~\ref{sec:assumptions}.
% First, the maximum-entropy
% method % (based on the \emph{Shannon} relative
% % entropy, as opposed to other entropy functions such as Burg's
% % \citeyear{burg1975})
% rests on some implicit assumptions about the probabilities for the long-run
% frequency distribution of activities, besides the assumptions just
% mentioned at the end of the previous paragraph. So the often-heard
% statement that it gives \enquote{the maximally unbiased (or non-committal)
%   distribution} must be taken with a grain of salt.
% \citep{jaynes1986d_r1996,portamana2009,portamana2017}. Second,
A maximum-entropy distribution like $\px(\yA \| M, N)$ is equivalent to the
zeroth-order Laplace approximation
\citep[\chap~4]{debruijn1958_r1961}{tierneyetal1986,strawderman2000} of
four distinct distributions for the larger population, which differ
numerically from one another in higher-order approximations:
\begin{enumerate}[label=(\roman*)]
\item\label{item:prob_freq} the most probable \emph{frequency} distribution
  for the total activity across the \emph{recorded} bins, %$\yFF(\yA)$,
\item\label{item:prob_binin} the \emph{belief} distribution for the value
  of the total activity at any time bin among the \emph{recorded} ones,
  %$\yp(\yA_{t} \| \yH)$ with $1\le t \le T$,
\item\label{item:prob_longrun} the most probable \emph{frequency}
  distribution for the total activity in a very long run of \emph{new} time
  bins, %$\ynuu(\yA)$,
\item\label{item:prob_binout} the \emph{belief} distribution for the value
  of the total activity at a \emph{new} time bin.
  % $\yp(\yA_{t} \| \yH)$ with $t<1$ or $t>T$
\end{enumerate}
Their mathematical relationship is presented more explicitly in
appendix~\ref{sec:inference_prob}. The
distribution~\eqref{eq:app_maxent_pop_first} can therefore be interpreted
in each of the four ways above, in this approximate sense. In the present
case the validity of the approximation decreases as the ratio $nN/T$
increases. The most accurate interpretation is~\ref{item:prob_freq}, and that
is how we prefer to interpret the distribution $\px$: as \emph{the most
  probable relative-frequency distribution} of the activity $\yA$. See
\sect~\ref{sec:assumptions} for further discussion about assumptions and
approximations.


\section{Example application: two data sets}
\label{sec:application}

We apply the model just described to two data sets publicly available in
the literature:
\begin{itemize}[wide]
\item The \textbf{first data set}, from \textcite[rat
  14147]{stensolaetal2012}, consists of $n=65$ neurons from rat Medial
  Entorhinal Cortex, simultaneously recorded for about 20 minutes. Their
  spikes are binned into $T=417\,641$ bins of $3\;\textrm{ms}$ width.
\item The \textbf{second data set}, from \textcite[supplementary
  material]{rostamietal2016_r2017}, %\texttt{data\_RS.npy}
  consists of $n=159$ neurons from macaque Motor Cortex, simultaneously
  recorded for about 15 minutes. Their spikes are binned into $T=300\,394$
  bins of $3\;\textrm{ms}$ width.
\end{itemize}

% In the followingFrom a Bayesian point of view, the maximum-entropy distribution $\px(\yA)$
% can be considered as \emph{the most probable frequency distribution} for
% the larger-population activity during the recording. But it also approximates
% three other distributions, as discussed in the previous section.

\bigskip

The maximum-entropy distribution for the larger population is calculated
using five moments, because this number provides almost as much information
as the full frequency distribution of the sample, as discussed in
\sect~\ref{sec:marginalization}. Figure~\ref{fig:plots_5mom} shows the
resulting distributions for three example values of
larger-population sizes: $N=1\,000$ (\textcolor{mygreen}{green diamonds}),
$N=5\,000$ (\textcolor{myred}{red circles}), $N=10\,000$
(small \textcolor{myblue}{blue squares}). The frequency density of the sample
activity is also shown (\textcolor{black}{black triangles}), and in the
plot it would be indistinguishable from the maximum-entropy density for
$N=n$, that is, applied at the sample level. We discuss the case of unknown
$N$ in \sect~\ref{sec:N}.
\begin{figure}[!p]
  \centering
\subcaption[]{\label{fig:plots_5mom_a}}%  
\includegraphics[width=\linewidth]{_scripts/newpaper_stensola_5mom_distributionsN_priorc.pdf}%
\\[1em]%\mbox{}\vfill\mbox{}%
\subcaption[]{\label{fig:plots_5mom_b}}%  
\includegraphics[width=\linewidth]{_scripts/newpaper_riehle_5mom_distributionsN_priorc.pdf}%
\\[1em]%\mbox{}\vfill\mbox{}%
\caption{Maximum-entropy distributions for the total activity of the larger
  population, using $M={}$five moments and assuming several population sizes $N$.
  In order to compare different $N$, the distributions are multiplied by
  $N$ (obtaining densities) and the activity divided by $N$ (obtaining a
  population average). The distributions for large $N$ have very different
  features from the measured frequency distribution, and how their peaks
  become more pronounced as $N$ increases.}
\label{fig:plots_5mom}
\end{figure}
% maxent_solutions_different_moments_prior2_6moments.nb

The most expensive calculation, for $N=10\,000$, takes less than 15 minutes
on a laptop with two $2\;\textrm{GHz}$ cores. In all cases the moments were
recovered with relative errors smaller than $10^{-12}$.
% 1000:  3*10^-13 \%
% 5000: 7*10^-13 \%, 581 s
% 10000: 6*10^-11 \%, 737 s

The figure shows that the distribution for the larger-population is more
peaked than the measured frequency distribution for the sample; their
difference increases with $N$. Most remarkably, for the first data set,
\fig~\ref{fig:plots_5mom_a}, the distribution for the larger population has
two distinct low-activity modes. For the second data set,
\fig~\ref{fig:plots_5mom_b}, the distribution presents a small shoulder
(barely visible) on the right of its mode.
% , suggestive of two activity regimes. We frame no hypotheses about the
% biological cause of these two modes (they could stem from the presence of
% different kinds of cells or modules)
Such features are clearly not present in the sample frequencies or in the
maximum-entropy distribution at the sample level. The model thus suggests
interesting features of the larger population. \textcolor{white}{If you
  find this you can claim a postcard from us.}

% maxent_solutions_different_moments_partial.nb
Here is a toy example of possible uses of this model, based on the first
data set. We could be interested in the hypothesis that two distinct cell
types or assemblies be present in the region where the recording was made.
Finding a larger-population distribution with two peaks, as in
\fig~\ref{fig:plots_5mom_a}, would provide some evidence for this
hypothesis. We might have reasons for suspecting that a specific set of the
sampled neurons is of the first type, and the remaining of the second type.
In the case of the first data set, indeed, 27 of the 65 sampled cells were
identified as grid cells belonging to 3--4 functional modules, the
remaining 38 \enquote{non-grid} cells being interneurons and unidentified
cells \citep{stensolaetal2012}. Could the two peaks in the distribution of
\fig~\ref{fig:plots_5mom} reflect the activities of grid versus non-grid
cells? We apply the model to these two sets of neurons individually, using
$\yNg=4\,150$ for the grid set and $\yNng=5\,850$ for the non-grid set, to
reflect their proportions (27/65 and 38/65) in the recorded sample. The
results are shown in \fig~\ref{fig:plots_cellsubsets_a}.
\begin{figure}[!p]
\centering
\subcaption[]{\label{fig:plots_cellsubsets_a}}%  
\includegraphics[width=\linewidth]{_scripts/newpaper3_gridvsnogrid.pdf}%
\\[1em]%
\subcaption[]{\label{fig:plots_cellsubsets_b}}%  
\includegraphics[width=\linewidth]{_scripts/newpaper3_gridconv.pdf}%
\\[1em]%
\caption{\subcaptionref{fig:plots_cellsubsets_a} Maximum-entropy
  distributions for larger subpopulations or grid and non-grid cells.
  \subcaptionref{fig:plots_cellsubsets_b} The convolution of the two
  subpopulation distributions is very different from the distribution for
  all cells; this means that the two subpopulation distributions are not
  independent.}
\label{fig:plots_cellsubsets}
\end{figure}
The distribution for the larger population of grid cells
(\textcolor{mygreen}{green triangles}) seems to have one broad peak, close
to the first peak of the larger-population distribution. The distribution
for the larger population of non-grid cells (\textcolor{myred}{red
  circles}) has two peaks instead, roughly at the same population-averaged
activities as the peaks of the larger-population distribution
(small \textcolor{myblue}{blue squares}) but closer in height. It would thus seem
that the population of grid cells is contributing to the first mode of the
larger population, but it is not its sole contributor. We can also assess
if the distributions for the two sets of neurons are independent. If they
were independent, the larger-population distribution would be given by
their convolution:
\begin{equation}
  \label{eq:convolution}
  P_{\text{full}}(\yA) = \sum_{\yA'} P_{\text{grid}}(\yA')\;
  P_{\text{non-grid}}(\yA-\yA') %, \qquad \yA \in \set{1,\dotsc,10\,000}
\end{equation}
where the index $\yA'$ runs from $\max(0,\yA-\yNng)$ to $\min(\yA,\yNg)$.
But this is not the case: 
figure~\ref{fig:plots_cellsubsets_b} shows that such convolution
(\textcolor{black}{black diamonds}) is quite different from the
larger-population distribution (small \textcolor{myblue}{blue squares}): the two
peaks of the former are closer in position and height than those of the
latter. The population distributions of grid and non-grid cells are
therefore \emph{not} independent: knowledge of the activity of either set
gives us some information about the activity of the other.

\medskip

The toy analysis above should not be taken literally, but just as an
illustration of the model's possible applications. Instead of a grid cell
\vs\ non-grid cell divisions we could for example have examined the modules
of the cells. In the absence of such auxiliary information about the cells,
one could in principle try all possible partitions of the recorded cells
into two groups, and check whether some partition leads to independent
distributions (computationally very expensive in practice); and so on.

The important point is that this model is computationally very cheap and
yet it can provide useful insights, even if just qualitative ones, and even
suggest new hypotheses.




\section{Quantifying the importance of higher-order correlations:\\
  the limitations of models at the sample level}
\label{sec:marginalization}

As mentioned in the Introduction, in the neurosciences the maximum-entropy
method has also been used to quantify the \enquote{cooperativity}
\citep[\eg][]{gersteinetal1985} or \enquote{interaction}
\citep[\eg][]{martignonetal1995,schneidmanetal2006,shlensetal2006} or
\enquote{synchrony} \citep[\eg]{bohteetal2000,amarietal2003} of neuronal
activity. Most, if not all, such applications use the method \emph{at the
  sample level}, that is, they find a distribution for the sample without
considering that it comes from a larger population. In this section we
discuss how our proposed application to a larger population bears on this
kind of quantification and compares with the application at the sample
level. %We also highlight some methodological caveats.% For concreteness we use an
% example with specific statistics: first and second moments (equivalent to
% mean and correlations), versus first-to-fourth moments. But our
% discussion holds for any sets of statistics.

\bigskip

As a first step, terms like \enquote{cooperativity}, \enquote{interaction},
and similar need to be translated into a more precise, operational
\citep[\cf][\chap~I]{bridgman1927_r1958} notion. Here we propose to base
this kind of analysis on the notion
of \emph{informational sufficiency}, because
\begin{enumerate*}[label=(\roman*)]
\item it seems closely related to those terms as used in the
  literature; \item it is mathematically connected with maximum-entropy
  distributions; \item it reveals important consequences of sampling.
\end{enumerate*}
Its meaning will now be briefly explained. For an extensive discussion of
this notion, which was probably first suggested by Kullback \amp\ Leibler
\parencite*{kullbacketal1951}, see Bernardo \etal\
\parencite*[\sect~4.5]{bernardoetal1994} and Jaynes \parencite*[\chap~8 and
\sect~14.2]{jaynes1994_r2003}.

Our belief distribution about the frequencies of the neuronal activities,
or about the activity in a new time bin, is conditional on all experimental
data and statistics we have collected. This belief distribution, however,
can be such that it remains unchanged if we discard part of the data or
statistics -- for example, the third- and higher-order empirical
moments. % --
% but it does change if we discard more.
This means that the discarded statistics are \emph{informationally
  irrelevant}: knowing them does not change our beliefs. The remaining
statistics -- in the example, first and second empirical moments -- are
\emph{informationally sufficient}.%  The qualification \enquote{minimally}
% will be omitted for brevity.
% This informational quasi-sufficiency is interesting because
% it may give clues about population organization and dynamics.

It seems to us that the informational sufficiency of a set of $m$ moments
captures the idea of the \enquote{importance} of the almost-simultaneous
activity of $m$ or fewer neurons: such simultaneous activity is
\enquote{unimportant} if its knowledge does not alter our forecasts. We are
thus speaking of informational importance, which may or may not imply or be
implied by biophysical importance -- \enquote{correlation is no proof of
  causation} \citep[\eg][]{simon1954,freedman2006}. But it lends itself to a
precise probabilistic analysis, with consequences determined by
mathematical theorems, as we shall now see.

Our question is whether some degree of \enquote{cooperation} is more
\enquote{important} than some other degree, or not. In terms of
informational sufficiency, its definition might suggest to compare the
belief distributions conditional on two different subsets of statistics.
But which measure should be used for such a comparison? and how to motivate
it? We prefer to take a different approach, leading to an automatic answer:
let us turn the informational sufficiencies of two different subsets of
statistics into two \emph{hypotheses}. We can then ask: which hypothesis is
more probable, given the data? Equivalently, what is the ratio of their
probabilities? By translating our question about \enquote{cooperativity}
into a hypothesis-testing problem about sufficiency, we can find a
well-defined answer determined by the probability-calculus. % The ratio of
% the two probabilities gives us a natural measure of the
% For example, do
% the data say that the first two empirical moments are probably enough to
% determine our beliefs? or is it more probable that the first four are
% needed instead?
We shall see that this approach is indeed more powerful than a simple
comparison of distributions.

But such a translation does more: through a mathematical theorem \citep[and
references therein]{andersen1970}[this theorem was first proved for
continuous probability distributions
by][]{koopman1936,pitman1936,darmois1935}[see][\sect~4.5.3 for further
references]{bernardoetal1994}, it leads to the automatic appearance of
maximum-entropy distributions: \emph{if a probability distribution for a
  repetitive phenomenon has some sufficient statistics, then it is a
  mixture of distributions of maximum-entropy form having those statistics
  as constraints, and vice versa}. 
% Thus, maximum-entropy distributions automatically appear (as
% approximations) from the probability-calculus when we test hypotheses
% about informational sufficiency.
The most important consequence of this theorem is that, by Bayes's formula,
the probability that some statistics are informationally sufficient can
approximately be found by first building the maximum-entropy distribution
constrained by those statistics, and then calculating the probability for
the whole data under that distribution. We shall shortly illustrate this
procedure with the data sets of \sect~\ref{sec:application}.

% For probabilities built with the maximum-entropy method, this means
% comparing those constrained on the subset and on the full set. This
% procedure is actually equivalent to comparing the probabilities of the two
% hypotheses about sufficiency, conditional on the \emph{full} data, assuming
% their pre-data probabilities to be equal. We show this equivalence below
% and perform the calculations for our first data set. \medskip

\medskip

Owing to the theorem above, studies about cooperativity (and similar ideas)
that use maximum-entropy can be naturally reinterpreted as studies about
informational sufficiency. An important fact must then be taken into
account: \emph{informational sufficiency is not preserved under sampling}.
If a probability distribution has some sufficient statistics, then its
sample marginals \emph{cannot} have the same sufficient statistics, and
vice versa \citep[\sect~3.1]{portamanaetal2015}; except for trivial cases
(\eg\ uniform probability distributions). % This generally holds also among
% different marginals: if a statistics is sufficient for some marginal, then
% it will not be sufficient for another marginal.
This impossibility is known in another guise in statistical mechanics: if a
system is described by a Gibbs state, its subsystems cannot be perfectly
described by Gibbs states \citep[\eg][and references
therein]{maesetal1999}. % Mathematically this impossibility comes from the
% Pitman-Koopman theorem mentioned above, and translates into the general
% impossibility of solving a system of independent equations with more
% equations than unknowns .

This fact affects the interpretation of the results obtained from
maximum-entropy distributions. For example, if means and pairwise
correlations seem informationally sufficient for a sample from a brain
area, then they are probably not sufficient for the larger population of
neurons constituting that area, and vice versa. If we are interested in the
cooperativity of the neurons in a brain area, it is therefore unreliable to
use a maximum-entropy distribution constructed only for a sample
thereof. % \citep[contrast this with some literature, \eg][]{martignonetal1995,bohteetal2000,schneidmanetal2006,shlensetal2006,barreiroetal2010,ganmoretal2011,granotatedgietal2013}
The model presented here avoids this problem because it targets the larger
population, not the sample alone.

\bigskip


We can compare the informational sufficiency of a set of moments, say
$M''\defd\set{1,\dotsc,m''}$, with respect to another, say
$M'\defd \set{1,\dotsc,m'}$, by introducing the quantity $\ydi(M'',M')$
defined as follows:
\begin{enumerate}[label=(\roman*)]
\item From each maximum-entropy distribution $\px(\yA \| M,N)$ for the
  larger population, \eqn~\eqref{eq:app_maxent_pop_first}, built from each
  set of constraints $M=M',M''$, calculate the marginal distribution for
  the sample by \eqn~\eqref{eq:conditional_hypergeometric_t}:
\begin{equation}
  \label{eq:sample_marginal}
  \pxx(\ya \| M, N) = \sum_{\yA}\yG_{\ya \yA}\px(\yA \| M,N),
  \qquad M=M',M''.
\end{equation}
\item Calculate the relative entropies of the measured frequency
  distribution $\yf$ with respect to each sample marginal from the previous
  step, and multiply them by the number of time bins $T$:
  \begin{equation}
    \label{eq:relentropy_T}
    T\,\sh[\yf; \pxx(\ya \| M,N)] \defd
    T  \sum_{\ya} \yff_{\ya} \log\frac{\yff_{\ya}}{\pxx(\ya \| M,N)},
    \qquad M=M',M''.
  \end{equation}
\item Take the difference: 
\begin{equation}
  \label{eq:measure_suffic}
  \begin{split}
  \ydi_{N}(M'',M') &\defd
  T \,\sh[\yf; \pxx(\ya \| M',N)]
  -
  T \,\sh[\yf; \pxx(\ya \| M'',N)]
  \\ &\equiv
T  \sum_{\ya} \yff_{\ya} \log\frac{
  \sum_{\yA}\yG_{\ya\yA}\px(\yA \| M'',N)}{
  \sum_{\yA}\yG_{\ya\yA}\px(\yA \| M',N)}.
\end{split}
\end{equation}
\end{enumerate}
The quantity $\ydi_{N}(M'',M')$ so defined is positive if $M''$ is more
probable to be informationally sufficient than $M'$, and negative
otherwise.

We show in appendix~\ref{sec:measure_suff} that $\ydi_{N}(M'',M')$ is
approximately equal to the log-ratio of the probabilities of the data $\yf$
conditional on the hypotheses $M''$ and $M'$:
\begin{equation}
  \label{eq:bayes_factor}
  \ydi_{N}(M'',M') \approx
  \log\frac{\pf(\yf \| M'',N)}{\pf(\yf \| M',N)}.
\end{equation}
This is their \emph{relative weight of evidence}
\citep[\chap~6]{good1950}{good1985}[\sect~1.4]{osteyeeetal1974}, the
logarithm of their \emph{relative Bayes factor}
\citep[p.~421]{jeffreys1936}[and references
therein]{kassetal1995}. %[\chaps~V, VI, A]{jeffreys1939_r1983}.
The exponential of $\ydi_{N}(M'',M')$ tells us how much more probable the
data $\yf$ are, conditional on the sufficiency of $M''$, than conditional
on the sufficiency of $M'$. We can also combine this measure with pre-data
probabilities for the two hypotheses to obtain the ratio of their
probabilities conditional on the data \citep[\cf][]{bretthorst2013}.


\bigskip

Let us illustrate the quantity just defined and the previous discussion
about sufficiency and sampling by using our data sets. We want to compare,
for example, three sets of constraints $M_{2}$, $M_{4}$, $M_{5}$,
consisting of the first two, four, five moments, for a larger population of
size $N=10\,000$. Figure~\ref{fig:plots_2_4mom} shows the resulting
distributions for the two data sets.
\begin{figure}[!p]
\centering
\subcaption[]{\label{fig:plots_2_4mom_a}}%  
\includegraphics[width=\linewidth]{_scripts/newpaper2_stensola_threemom_distributions10000_priorc.pdf}%
\\[1em]%
\subcaption[]{\label{fig:plots_2_4mom_b}}%  
\includegraphics[width=\linewidth]{_scripts/newpaper2_riehle_threemom_distributions10000_prior2.pdf}%
\\[1em]%
\caption{Maximum-entropy distributions for the total activity of the larger
  population, based on the first two, four, and five moments, and a
  larger-population size $N=10\,000$. The two- and four-moment
  distributions are very different, suggesting that two moments are not
  informationally sufficient. The four- and five-moment distributions are
  almost indistinguishable, suggesting that four and five moments are
  equally likely to be sufficient.}
\label{fig:plots_2_4mom}
\end{figure}
% maxent_solutions_different_moments_prior2_6moments.nb
There is a clear visual difference between the distributions based on two
(\textcolor{myred}{dashed red}) and on four (\textcolor{myblue}{solid
  blue}) moments; for the first data set, panel
\subcaptionref{fig:plots_2_4mom_a}, they even differ in the number of
modes. On the other hand there is barely any visual difference between the
distributions based on four and on five (\textcolor{myyellow}{dotted
  yellow}) moments. From these visual differences we immediately conclude
that two moments are less likely to be informationally sufficient than
four; but four should be almost as likely as five. We shall see that the
measure~\eqref{eq:measure_suffic} reflects this swift visual judgement.

Compare these plots with those obtained applying the model at the sample
level (equivalent to $N=n$) to the first data set, shown in
\fig~\ref{fig:comparisons_marginals}. The distributions based on two
(\textcolor{myred}{red circles}), four (\textcolor{myblue}{blue empty
  squares}), and five (\textcolor{myyellow}{yellow lozenges}) moments can
only be distinguished on a logarithmic scale, and only in their tails. The
measured frequencies (black triangles) are also shown. Two moments seem
slightly less likely to be sufficient than four; and four and five moment
seem about equally likely. But it is difficult to visually make these
comparisons.

\medskip

% NUMERICAL RESULTS
%
% relentropies of sample2mom, sample4mom, sample5mom, marg2mom, marg4mom, marg5mom
% -322.72,-42.8005,-39.4917,-85.7626,-4.23725,-4.20059 nats =
% -1401.56,-185.88,-171.51,-372.462,-18.4021,-18.2429 dB
%
% difference in relentropies between 4 mom and 2 mom:
%** sample level: 279.92 nats = 403.839 bit = 1215.68 dB
%** marg level: 81.525 nats =  117.616 bit = 354.06 dB
% difference in relentropies between 5 mom and 4 mom:
%** sample level: 3.30882 nats = 4.77362 bit = 14.37 dB
%** marg level: 0.0366623 nats = 0.0543836 bit = 0.159223 dB
% correspond to prob ratios of
% sample: 3.69524\times 10^{121}, 27.3529
% marg: 2.5469\times 10^{35}, 1.03734
% 
% comparison N vs n:
% 2 mom: 236.958 nats = 341.858 bit = 1029.09 dB, 8.11758\times 10^{102}
% 4 mom: 38.5632 nats = 55.635 bit = 167.478 dB, 5.59496*10^16
% 5 mom: 35.2911 nats = 50.9157 bit = 153.267 dB, 2.12185*10^15
%
% comparison_ME_sample_full_levels.nb
%
Calculation of the weight of evidence~\eqref{eq:measure_suffic} for the
first data set and $N=10\,000$, corresponding to
\fig~\ref{fig:plots_2_4mom_a}, leads to the following results:
\begin{equation}
  \label{eq:diff_suff_N}
  \begin{split}
  \ydi_{N}(M_{4},M_{2}) &= 118\;\bit \equiv 35.4\;\hart,
\\    \ydi_{N}(M_{5},M_{4}) &= 0.0544\;\bit \equiv 0.0159\;\hart,
  \end{split}
\end{equation}
where the Hartley ($\hart$) denotes base-10 logarithms
\citep[\sect~C.4]{iso2009}[called \enquote{ban} and used by Turing and Good
in their code-breaking work at Bletchley Park:
see][]{good1985%,good1950,good1969
}[\sect~4.2]{jaynes1994_r2003}. In words, our data are 35 orders of
magnitude more probable by assuming sufficiency of $M_{4}$ than by assuming
sufficiency of $M_{2}$. But the data are about as probable
($10^{0.016} = 1.04$) by assuming sufficiency of $M_{4}$ as of $M_{5}$. The
qualitative visual information of \fig~\ref{fig:plots_2_4mom_a} was
therefore reliable.

For the model at the sample level, that is, with $N=n$, the weights of
evidence are
\begin{equation}
  \label{eq:diff_suff_n}
  \begin{split}
  \ydi_{n}(M_{4},M_{2}) &= 404\;\bit \equiv 122\;\hart,
\\    \ydi_{n}(M_{5},M_{4}) &= 4.77\;\bit \equiv 1.44\;\hart.
\end{split}
\end{equation}
So the application at the sample level says that the data are 120 orders of
magnitude more probable conditional on $M_{4}$ than on $M_{2}$, and
$10^{1.4}=25$ times more probable conditional on $M_{5}$ than on $M_{4}$.
These results are different from those of the application at the larger
population, and they would be hard to deduce visually from the
distributions of \fig~\ref{fig:comparisons_marginals}.



% The conclusions about \enquote{cooperativity} or \enquote{interaction} that
% we reach by applying the maximum-entropy method to the larger population,
% as proposed here, are therefore different from those applying it at the
% sample level -- and also visually clearer. The plot in 
% \fig~\ref{fig:plots_2_4mom_a}, showing the two clearly different
% distributions constrained by two and four moments, should be compared with
% the plot for the distributions obtained at the sample level under the same
% constraints, shown in \fig~\ref{fig:comparisons_marginals} as
% \textcolor{myred}{red circles} and \textcolor{myblue}{blue squares}. The
% measured frequencies (black triangles) are also shown. It is necessary to
% use a logarithmic scale to see the differences, which appear mainly in the
% tail.
\begin{figure}[!p]
\centering
% \includegraphics[width=\linewidth]{_scripts/comparison_full.pdf}%
% \\[3em]%
\includegraphics[width=0.92\linewidth]{_scripts/trunct2_comparison_marginals.pdf}%
\caption{Maximum-entropy distributions constructed at the sample level,
  using the first two, four, and five moments; and sample marginal of the
  distribution constructed with five moments for a larger population with
  $N=10\,000$ (dotted yellow in \ref{fig:plots_2_4mom_a}). The
  distributions are multiplied by $T$, obtaining absolute frequencies; note
  that activities with frequency lower than 1 cannot be observed in $T$
  bins. The distributions at the sample level are only distinguishable on a
  logarithmic scale. The five-moment distribution constructed at the
  larger-population level fits the frequency data better than the one at
  the sample level.}
\label{fig:comparisons_marginals}
\end{figure}
% comparison_ME_sample_full_levels.nb

\medskip

% 4 mom: 38.5632 nats = 167.478 dB, 5.59496*10^16
% 5 mom: 35.2911 nats = 153.267 dB, 2.12185*10^15

It is possible to use the measure~\eqref{eq:measure_suffic} and its
probabilistic interpretation~\eqref{eq:bayes_factor} to compare the
probability of the data $\yf$ conditional on the hypothesis $(M_{5},N)$ of
five-moment sufficiency at the larger-population level, $N=10\,000$, versus
the hypothesis $(M_{5}, n)$ of five-moment sufficiency at the sample level,
$N=n=65$. We obtain
\begin{equation}
  \label{eq:diff_suff_n}
  \ydi[(M_{5}, N), (M_{5}, n)] = 50.9\;\bit \equiv 15.3\;\hart.
\end{equation}
% that is, the data are 15 orders of magnitude more probable under the first
% hypothesis than under the second. 
This result is also shown in
\fig~\ref{fig:comparisons_marginals}. We see that the five-moment
distribution constructed at the larger-population level
(\textcolor{myyellow}{yellow stars}) fits the measured frequency
distribution (black triangles) more closely than the corresponding
sample-level distribution (\textcolor{myyellow}{yellow lozenges}).



\bigskip

% We close this section with some references. The notion of informational or
% predictive sufficiency, probably first presented by Kullback \amp\ Leibler
% \parencite*{kullbacketal1951}, is discussed in Bernardo \etal\
% \parencite*[\sect~4.5]{bernardoetal1994}, who explain its connections with
% Fisher's\parencite*{fisher1922} parametric sufficiency. Jaynes
% \parencite*[\chap~8 and \sect~14.2] shows the connections with information
% theory. The mathematical relation between sufficiency and the exponential
% family comes from the Pitman-Koopman
% theorem \citep{koopman1936,pitman1936,darmois1935} for continuous variables.


% \mynote{Short reasoned bibliography about sufficient statistics and its
%   relation to maximum-entropy?} blah blah \footnote{For more technical results and
%   connections with the notion of symmetry see \eg\
%   \cite{darmois1935,neyman1935,koopman1936,pitman1936,halmosetal1949,bahadur1954,berk1972,lauritzen1974,lauritzen1982_r1988,lauritzen2007,cifarellietal1980,cifarellietal1981,diaconisetal1981,diaconis1992,furmanczyketal1998,fortinietal2000,nogalesetal2000,kallenberg2005,ayetal2015}.}
% This mathematical result comes from the Pitman-Koopman
% theorem\citep{kullbacketal1951}[the
% notion goes back to][]{fisher1922} \citep{jaynes1982b}[\sect~4.5.4]{bernardoetal1994}
% \citep{koopman1936,pitman1936,darmois1935}[for later analyses and the
% discrete case
% see][]{andersen1970,%discrete
% denny1967,denny1972}




% But does the application at the larger-population level lead to appreciably
% different results from that at the sample level? after all we're interested
% in an approximate informational sufficiency, not in an analytically exact
% one. A correct answer can only be given case by case.
% Figure~\ref{fig:comparisons_marginals} gives a graphical answer to this
% question in the case of our first data set. The upper plot shows the
% maximum-entropy distributions for a larger population of $10\,000$ neurons,
% constructed from two moments (\textcolor{mygreen}{\textbf{--~--~--}}) and
% from four moments (\textcolor{mypurpleblue}{\textbf{------}}). The lower
% plot shows the maximum-entropy distributions for the sample, from two
% moments (\textcolor{mygreen}{$\bm{\blacktriangledown}$}) and from four moments
% (\textcolor{mypurpleblue}{$\bm{\square}$}), and the distributions obtained
% by marginalizing the larger-population distribution to the sample size, from two
% moments (\textcolor{mygreen}{$\bm{\blacktriangledown}$}) and from four
% moments (\textcolor{mypurpleblue}{$\bm{\blacksquare}$}). The measured
% frequency distribution (\textcolor{myred}{\LARGE$\bm{\bullet}$}) is also
% shown. We note the following:
% \begin{itemize}
% \item The application to the larger-population (upper plot) leads to two
%   completely different distributions (even different number of modes);
%   clearly the two sets of statistics are not even approximately equivalent
%   for inferential purposes.
% \item The application at the sample level leads to deceivingly similar
%   distributions (\textcolor{mygreen}{$\bm{\blacktriangledown}$} and
%   \textcolor{mypurpleblue}{$\bm{\square}$}), which can only be distinguished on a
%   logarithmic scale. This could lead to the erroneous conclusion that the
%   two statistics are approximately equivalent.
% \item The difference between marginals from the larger-population distribution
%   and the distributions obtained at the sample level (filled vs empty
%   markers) is larger than the difference between different statistics
%   (triangles vs squares).
% \item The marginals of the application to the larger population are closer to
%   the measured frequencies than the distributions obtained at the sample
%   level (although this closeness is not a valid criterion for their
%   goodness).
% \end{itemize}
% Therefore, the maximum-entropy application at the sample level and at the
% larger-population level lead to different results; the latter is not only more
% meaningful but also superior because it shows clearly the different
% informational sufficiency of the two sets of statistics.


\section{The size $N$ of the larger population}
\label{sec:N}

% results for post-data probabilities of N given data and 5 moments:
% N=65,1000,2000,5000,10000
% 1.40169\times 10^{-16}, 0.0439602, 0.139657, 0.251064, 0.297726, 0.267594
%
% using 1/N pre-data prob:
% 1.04099\times 10^{-14}, 0.212211, 0.337085, 0.242394, 0.143722, 0.0645882
% 
% maxent_solutions_different_moments_prior2_6moments.nb

The calculations presented in the preceding sections depend on the size $N$
of the larger population, which of course cannot be the whole brain. How
large should or can $N$ be in our formulae?

The answer comes by examining the origin of our belief distribution for the
activity of the sample \emph{conditional} on the activity of the larger
population. It is given by the hypergeometric
distribution~\eqref{eq:conditional_hypergeometric}, reprinted here:
\begin{equation}
  \labelbis{eq:conditional_hypergeometric}
  \yG_{\ya\yA} \defd \pf(\ya \|\yA, n,N)=
    \binom{\yA}{\ya}\binom{N-\yA}{n-\ya}\binom{N}{n}^{-1}
\equiv \binom{n}{\ya}\binom{N-n}{\yA-\ya}\binom{N}{\yA}^{-1}.
\end{equation}
This distribution leads to the equality of factorial
moments~\eqref{eq:equal_m_factorialmoments}, on which our model rests. This
conditional probability, characteristic of \enquote{drawing without
  replacement}
\citep[\eg][\chap~3]{jaynes1994_r2003}, % [\sect~4.8.3]{ross1976_r2010}[\sect~II.6]{feller1950_r1968}
ensues when we believe, with the same degree, that any of the $N$ units
(balls or neurons) could be drawn or sampled. Thus, consider the pool of
all neurons which we believe -- with an \emph{equal} degree of belief for
each neuron -- could have been recorded. Then $N$ is the size of that pool.
$N$ therefore depends on factors such as: the shape, dimensions, and
technical specifications of the recording probe; the inaccuracy in the
insertion of the probe, leading for example to slightly different insertion
points or angles; the density of neurons around the probe. But it also
depends on the homogeneity of the brain region where the recording is made:
if we believe that it does not matter whether the probe had been inserted
in some other point of the same brain region, then
formula~\eqref{eq:conditional_hypergeometric} is appropriate. $N$ can
therefore only be assessed case by case.

Whether this equality in our degrees of belief is justified, or for which
sampling procedures it can be justified, is the fundamental, deep question
of sampling theory, which we cannot discuss here \citep[see \eg\ the
discussions and references
in][]{ericson1969,smith1976}. % {kruskaletal1979c,kruskaletal1980}
It is of course possible to probabilistically assess, case by case, whether
this symmetry assumption is appropriate. But we must also be aware that
such assessment in turn rests on analogous symmetry assumptions at a higher
level, and so on. The probability-calculus, just like the truth-calculus,
cannot yield conclusions without premisses
\citep{hailperin2011}[p.~182]{johnson1924}.

The plots of \fig~\ref{fig:plots_5mom} nevertheless suggest that the order
of magnitude of $N$ ought to be enough for qualitative inferences. As
discussed in the next section, even if the exact number $N$ were known, the
maximum-entropy distribution ought to be interpreted qualitatively or
semi-quantitatively.

\medskip

If our uncertainty about $N$ spans several orders of magnitude we can
assign a probability to the possible values of $N$ based on our background
information and the frequency data $\yf$, similarly to the procedure in
\sect~\ref{sec:marginalization} and appendix~\ref{sec:measure_suff}. The
probability for the data $\yf$ conditional on a set of constraints $M$ and
size $N$ is the exponential of the negative relative
entropy~\eqref{eq:relentropy_T}:
\begin{equation}
  \label{eq:prob_data_from_N}
  \pf(\yf \| M, N) \approx \binom{T}{T\yf}\;\prod_{a}\pxx(\ya \| M, N)^{T\,\yff_{a}}
  \approx \exp\bigl\{  -T\,\sh[\yf; \pxx(\ya \| M, N)]\bigr\}.
\end{equation}
If our pre-data probability for $N$ is $p(N)$, then by Bayes's theorem
\begin{equation}
  \label{eq:prob_N_from_data}
  \pf(N \| M, \yf) \propto
  p(N)\;\exp\bigl\{-T\,\sh[\yf; \pxx(\ya \| M, N)]\bigr\}.
\end{equation}

Here is an illustrative example with our first data set. Suppose our
uncertainty spans slightly more than one order of magnitude, from $N=1\,000$
to $N=20\,000$. Divide this range roughly into thirds of order of
magnitude, considering the values
$N\in \set{10^{3}=1\,000, 10^{3+1/3}\approx 2\,000, 10^{3+2/3}\approx
  5\,000, 10^{4}=10\,000, 10^{4+1/3}\approx 20\,000}$. Assuming the set $M$
of first five moments to be sufficient, formula~\eqref{eq:prob_data_from_N}
gives
\begin{equation}
  \label{eq:data1_f_from_N}
  \begin{gathered}
  \pf(\yf \| N=1\,000, M) =  0.00222,\qquad
  \pf(\yf \| N=2\,000, M) =  0.00704,\\
  \pf(\yf \| N=5\,000, M) =  0.0127,\qquad
  \pf(\yf \| N=10\,000, M) =  0.0150,\\
  \pf(\yf \| N=20\,000, M) =  0.0135.
\end{gathered}
\end{equation}
Since our uncertainty regards a scale factor, it can be represented by
equal pre-data probabilities of $1/5$ about these partial orders of
magnitude. Thus from~\eqref{eq:data1_f_from_N} we find
\begin{equation}
  \label{eq:data1_N_from_f}
  \begin{gathered}
  \pf(N=1\,000 \| \yf, M) =  0.044,\qquad
  \pf(N=2\,000 \| \yf, M) =  0.140,\\
  \pf(N=5\,000 \| \yf, M) =  0.251,\qquad
  \pf(N=10\,000 \| \yf, M) =  0.298,\\
  \pf(N=20\,000 \| \yf, M) =  0.267,
\end{gathered}
\end{equation}
which gives a slightly higher probability to $N=10\,000$.

\bigskip

In this way we can make inferences % -- for example about the sufficiency of
% a set of moments, or about the marginal sample distribution --
that take into account our uncertainty about $N$. We must make sure to
avoid circularities: for example, we cannot assume a set of moments
$M$ to be sufficient and assess our uncertainty about $N$ conditional on
it, and then use this uncertainty to assess the sufficiency of $M$. But we
can approximately assess the sufficiency of a subset of moments
$M' \subset M$ taking into account the uncertainty of $N$ conditional on
$M$.

A rigorous assessment would involve a more expensive, full Bayesian
calculation; but such calculation would make the maximum-entropy
method superfluous. We discuss this final point in the next section.


\section{Assumptions and approximations that define the model}
\label{sec:assumptions}

The possible uses of the model here presented depend on the assumptions and
approximations which define it: three main assumptions and two main
approximations all in all. They become clear within a full-fledged
probabilistic approach. We summarize them here, leaving a more mathematical
sketch to appendix~\ref{sec:inference_prob}.

The first and most important assumption, typical of this kind of
maximum-entropy application in the neurosciences, can be expressed in two
equivalent ways:
\begin{enumerate*}[(\roman*)]
  % \item the order of the time bins has no relevance;
\item the frequencies of the activity during a recording are
  informationally sufficient for inferences about unrecorded times; \item
  our degree of belief about the sequence of activities is invariant under
  permutations of the time bins.
\end{enumerate*}
Equating time averages and expectations, \eqn~\eqref{eq:final_constraint},
would not make sense without this assumption. Its mathematical consequence
is to relate the probability for the frequency distribution $\yF$ for the
activity of the larger population during the recording, to the probability
for the \emph{long-run} frequency distribution. In statistical physics, as
a comparison, the maximum-entropy method is instead used with the
assumption of permutation-invariance across \emph{repetitions} of the same
experiment, because its constraints represent macroscopically
\emph{reproducible} quantities and kinematics, even in non-equilibrium
\citep{jaynes1985b_r1996,%jaynes1957,%tikochinskyetal1984,
  grandy1988,%berryetal1988,
  deroecketal2005_r2006%,touchette2009
}[for experiments \eg][]{guntonetal1983}. The kinetic-Ising models
mentioned in \sect~\ref{sec:intro} are closer to this latter approach.

The second assumption, typical of maximum-entropy in general, concerns the
probability distribution for the long-run frequencies ensuing from the
first assumption \citep[and references therein]{portamana2017}. It is taken
to be either an \enquote{entropic prior}
\citep{skillingetal1984,rodriguez1991,neumann2007}, which takes into
account the multiplicity of long-run frequencies, or a prior for a model
with sufficient statistics, which expresses the sufficiency of the moments
under study, according to the theorem discussed in
\sect~\ref{sec:marginalization}. This assumption is crucial for the
maximum-entropy method based on the Shannon entropy. Alternative prior
densities lead to alternative maximum-entropy methods; for example one
based on the Burg entropy \citep{burg1975} if a Dirichlet density is used
\citep{jaynes1986d_r1996}. The entropic prior depends on a reference
distribution and on a coefficient $L$ which determines the sharpness of the
prior's peak. The computations of \sects~\ref{sec:application}
and~\ref{sec:marginalization} were insensitive to the choice of the
reference distribution in their significant digits. We used three
biologically reasonable alternatives: a flat distribution, one linearly
decreasing with $\yA$, and a normal one with a broad peak at low
activities.

The third assumption is specific to our application: all subsets of $\yA$
neurons in the larger population are equally likely to give rise to total
activity $\yA$, at every time bin. A mathematically equivalent assumption
is that $n$ neurons are sampled anew at each time bin. This assumption
allows us to factorize the joint probability of the sequence of activities
and apply the sampling relation~\eqref{eq:conditional_hypergeometric} of
appendix~\ref{sec:derivation_maxent_fullnet} at each bin.

Finally, the maximum-entropy method is typically equivalent to
approximating the number $T$ of time bins, the prior coefficient $L$, and
the ratio $T/L$ to infinity (that is, $T\gg L \gg 1$). These approximations
lead to the identification of the four distributions listed at the end of
\sect~\ref{sec:model}: the recorded frequency distribution becomes equal to
the long-run one, and also to the probability distribution for the activity
at recorded and new time bins. In our application it also leads to the
equality of the hypergeometric
distribution~\eqref{eq:conditional_hypergeometric} and the \emph{conditional}
frequencies of the sample activity $\ya$ given $\yA$. The goodness of this
approximation decreases as the ratio -- roughly $nN/T$ -- between the
number of possible joint states and the number of bins increases, and also
as observed quantities approach their mathematical bounds, for example some
frequencies $\yff_{\ya}=0$. Without these approximations it is still
possible to interpret the maximum-entropy
distribution~\eqref{eq:app_maxent_pop_first} as the mode of the probability
distribution for the frequencies $\yF$, that is, the most probable
recorded frequency.


\medskip

The first and third assumptions above are the least realistic, but they
make sense as reference hypotheses for judging the informational relevance
of correlations of some order, as discussed in
\sect~\ref{sec:marginalization}, to be compared with hypotheses about more
realistic time dependences and
clustering. %They may turn out to approximate well more realistic assumptions.

The $T= \infty$ approximation (together with the third assumption above)
can break our model for specific distributions of sample frequencies, if
too many moments are used. For example, take $n=2$ with frequency
distribution $\yf=(0.5,0,0.5)$ for the activities $\ya=0,1,2$; that is, the
sample is silent half of the time and completely active half of the time.
The first two normalized factorial moments have both value $0.5$, and
together they uniquely determine $\yf$. If we now take $N=3$, it can be
proved that \eqn~\eqref{eq:app_maxent_pop_first}, constrained on those two
moments, has no solution. The reason is clear from a sampling perspective:
if there is zero probability of sampling one active neuron, then the larger
population cannot have any active neurons at all, and a fortiori there must
be zero probability of sampling two active ones. But this reasoning holds
only for each single time bin, not for the averages over time bins. The
$T=\infty$ approximation identifies the two instead. Outside of this
approximation there are no contradictions because the conditional
frequencies and conditional probabilities are kept distinct.

The two infinity approximations are at the edge of their validity range, in
our examples. For this reason we interpret the maximum-entropy distribution
as \emph{the most probable frequency distributions}, as mentioned at the
end of \sect~\ref{sec:model}. In high dimensions, however, the features of
the most plausible distribution may not be \emph{typical} of the majority
of plausible distributions. An important question, then, is \emph{which
  features of the maximum-entropy frequency distribution are typical of the
  majority of plausible frequency distributions?} -- especially for
features such as the number, position, and heights of modes; or tail
behaviour. Preliminary studies using a full probability calculation (in
preparation, to be published separately) indicate that several
semi-quantitative features of the maximum-entropy distribution are indeed
typical. For the first data set, for example, most frequency distributions
have a high-frequency regime at a population-averaged activity of around
0.075, as in \fig~\ref{fig:plots_5mom_a}, but their second high-frequency
regime is typically lower and appears for $N>5\,000$. Most frequency
distributions have two or three very close modes within each high-frequency
regime, rather than one. The same preliminary studies also indicate that
these typical features are insensitive to alternative prior choices in the
second assumption above, for example using a Dirichlet prior instead of an
entropic one.

This suggests that the maximum-entropy approximation presented here can be
favourably used to get a qualitative or semi-quantitative idea of what a
fully probabilistic calculation could yield -- but with a far lower
computational cost than the latter. For example, for $N=1\,000$ the
maximum-entropy distribution can be calculated in ten minutes on a laptop,
whereas the full probability calculation takes several days on a
high-performance computing cluster.



% \medskip

% The second point is that our degrees of belief about the frequency
% distribution for the larger population depend not only on the measured data in
% the sample, but also on our pre-data beliefs $\yI$ about the distribution.
% Which assumptions lead to the maximum-entropy result? This distribution
% appears when our initial belief about the possible frequency distributions
% $\yF$ is quantified by an entropic prior
% \citep{neumann2007,rodriguez1991,skilling1998,catichaetal2004,portamana2017}:
% \begin{equation}
%   \label{eq:entropic_prior}
%   \pf(\yF \| \yI) \propto  \exp[-\yL\;\re(\yF;\yR)]
%   \approx \binom{L}{L\yFF_{0}, \dotsc, L\yFF_{N}}\;
%   \prod_{\yA}{\yRR_{\yA}}^{L\yFF_{\yA}}
% \end{equation}
% where $\re(\yF;\yR)\defd \sum_{\yA}\yFF_{\yA}\log(\yFF_{\yA}/\yRR_{\yA})$ is
% the relative entropy or discrimination information
% \citep{kullback1987,jaynes1963,hobson1969,hobsonetal1973}, $\yR$ is the
% reference distribution, and $\yL$ a positive parameter. The approximate
% equality (obtained through Stirling's approximation), where the large
% parentheses denote a multinomial coefficient, shows that this prior belief
% is proportional to the number of ways in which the distribution $\yF$ can
% be realized in $L$ time bins. The parameter $L$ roughly quantifies how many
% time bins our data set must have to affect our initial belief. The
% maximum-entropy approximation is valid when $L$ is large, but small
% compared to the sharpness of the constraints on $\yF$; in our case this
% means $\yL \approx 10$, give or take an order of magnitude.



\section{Summary and discussion}
\label{sec:summary_discussion}

The reason why we record neurons from a brain region is to have an idea of
the activity of the other neurons in that region, similarly to survey
sampling. If we thought that \emph{precisely} the recorded ones were
exceptions or completely unrelated to the others surrounding the recording
probe, our recording would serve little purpose. Probabilistic models that
allow us to make direct inferences about the larger population are
therefore useful. But they become complex and sensitive if we ask for very
detailed inferences.

We have presented a model that allows us to make direct inferences about a
larger population from very small samples. Its inferences regard only the
frequencies of the total activity of the larger population, but such
limited inferential scope makes the model computationally very cheap and,
most important, it still leads to inferences that are not at all trivial
and could not be made just by looking at the sample and its distributions.
We showed this by applying the model to two real data sets.

By \enquote{inference} we mean the quantification of a degree of belief
about possible observations, based on specific assumptions; not the
response of an oracle. The model here presented gives us a \emph{forecast}
based on specific assumptions discussed in \sect~\ref{sec:assumptions}, and
therefore has two main uses, just like any other probabilistic model. If
our assumptions are only hypotheses under testing, then the subsequent
confirmation or confutation of the forecast will increase or decrease our
confidence in those hypotheses. If we strongly trust our assumptions, on
the other hand, then we rely on the forecast and even use it in place of
actual observations, especially if the latter are difficult to make. (In
fact, we typically shift from the first to the second use.) In the
preceding sections we have given examples of additional hypotheses that
could be explored with the present model. The assumptions of the model
itself, discussed in \sect~\ref{sec:assumptions}, can either be under test
or relied upon, depending on the case.

Some approximations implicit in the model make its numerical results
semi-quantitative with respect to an exact analysis.%  The results should not
% be off by more than a few percentages for the population-averaged
% activity and for the frequency density.
This imprecision, however, is compensated by the
extreme computational cheapness of the model, far less costly than an exact
analysis. The model is therefore very convenient at least for essaying
hypotheses. We hope to present an exact quantitative comparison with the
full probabilistic approach in a future publication.



% \mynote{cite  \cite[p.~440]{maxwell1873c}}

% We have presented a procedure to construct the most plausible frequency
% distribution of population-averaged activities of a population of neurons,
% given the recording about a small sample thereof. This procedure combines
% the maximum-entropy method and basic identities from sampling theory. From
% the application to two real data sets we saw that the frequency
% distributions obtained with our procedure can have features very different
% from the one measured in the sample, such as multiple modes. This procedure
% can also be used with moment constraints of different order -- means,
% population-averaged pairwise correlations, or higher-order correlations --
% thus giving an approximate assessment of the informational sufficiency of
% specific subsets of moments. In fact, we saw that the application of
% maximum-entropy only at the sample level leads to misleading results about
% this kind of sufficiency questions.




%%\setlength{\intextsep}{0.5ex}% with wrapfigure
%\begin{figure}[p!]%{r}{0.4\linewidth} % with wrapfigure
%  \centering\includegraphics[trim={12ex 0 18ex 0},clip,width=\linewidth]{maxent_saddle.png}\\
%\caption{caption}\label{fig:comparison_a5}
%\end{figure}% exp_family_maxent.nb


%%%%%%%%%%%%%%%%%%%%%%%%%%%%%%%%%%%%%%%%%%%%%%%%%%%%%%%%%%%%%%%%%%%%%%%%%%%%
%%% Acknowledgements
%%%%%%%%%%%%%%%%%%%%%%%%%%%%%%%%%%%%%%%%%%%%%%%%%%%%%%%%%%%%%%%%%%%%%%%%%%%% 
\begin{acknowledgements}
  This work is financially supported by the Kavli Foundation and the Centre
  of Excellence scheme of the  Research Council of Norway.
  
  We thank Iv\'an Davidovich for very insightful and valuable comments on a
  previous draft of this manuscript. PGLPM thanks the staff of the NTNU
  library for their always prompt support; Mari, Miri, Emma for
  continuous encouragement and affection; Buster Keaton and Saitama for
  filling life with awe and inspiration; the developers and maintainers of
  \LaTeX, Emacs, AUC\TeX, Open Science Framework, Python, Inkscape, Sci-Hub
  for making a free and unfiltered scientific exchange possible.
%\rotatebox{15}{P}\rotatebox{5}{I}\rotatebox{-10}{P}\rotatebox{10}{\reflectbox{P}}\rotatebox{-5}{O}.
%\sourceatright{\autanet}
\end{acknowledgements}

%%%%%%%%%%%%%%%%%%%%%%%%%%%%%%%%%%%%%%%%%%%%%%%%%%%%%%%%%%%%%%%%%%%%%%%%%%%%
%%% Appendices
%%%%%%%%%%%%%%%%%%%%%%%%%%%%%%%%%%%%%%%%%%%%%%%%%%%%%%%%%%%%%%%%%%%%%%%%%%%% 
%\clearpage
% %\renewcommand*{\appendixpagename}{Appendix}
% %\renewcommand*{\appendixname}{Appendix}
% %\appendixpage
\appendix

\section{Derivation of the maximum-entropy distribution}
\label{sec:derivation_maxent_fullnet}

Here is a summary derivation of the maximum-entropy distribution for the
larger population constrained by a set of factorial moments.%  For further details
% see Porta Mana \etal\ \citey{portamanaetal2015}.

\bigskip

First of all the sampling relation. We have a population of $N$ units
(neurons), $\yA$ of which have some specific property (being active in a
specific time bin), and we sample $n$ of the $N$ units in such a way
(partly uncontrollable experimental procedure) that we are led to believe,
with the same degree, that any of the $N$ units could have been sampled.
The probability that $\ya$ of the $n$ sampled neurons are active is then
given by the hypergeometric distribution
\begin{equation}
  \label{eq:conditional_hypergeometric}
  \yG_{\ya\yA} \defd \pf(\ya \|\yA, n,N)=
  \binom{\yA}{\ya}\binom{N-\yA}{n-\ya}\binom{N}{n}^{-1}
\equiv  \binom{n}{\ya}\binom{N-n}{\yA-\ya}\binom{N}{\yA}^{-1}
\end{equation}
typical of \enquote{drawing without replacement}
\citep[\eg][\chap~3]{jaynes1994_r2003}. %[\sect~4.8.3]{ross1976_r2010}[\sect~II.6]{feller1950_r1968}[\sects~2.1, 3.2]{jeffreys1939_r1983}.
In the following we leave $n$, $N$ implicit in the conditional. If we are
uncertain about the number $\yA$, with belief $P(\yA)$, then by the theorem
of total probability our belief about $\ya$ is
\begin{equation}
  \label{eq:belief_sample_unknown_A}
  p(\ya)= \sum_{\yA} \yG_{\ya\yA} \;P(\yA).
\end{equation}

From the definition of normalized factorial
moment~\eqref{eq:def_factorial_moment_a}, the expression for the
hypergeometric distribution~\eqref{eq:conditional_hypergeometric}, and the
relation~\eqref{eq:belief_sample_unknown_A} between our beliefs about $\ya$
and $\yA$, and using some combinatorial juggling
\citep%[\chaps~I--IV]{whitworth1867_r1965}[\chap~II]{feller1950_r1968}
[appendix~A and references therein]{portamanaetal2015}, one can prove the
equality~\eqref{eq:equal_m_factorialmoments} between the factorial moments
of sample and larger population.

\bigskip

For the construction of a maximum-entropy distribution from generic
expectation constraints see Jaynes
\citey{jaynes1963}[\chap~11]{jaynes1994_r2003}, and Hobson \etal\
\citey{hobsonetal1973} and Sivia \citey[\sect~5.2.2]{sivia1996_r2006} for
the minimum-relative-entropy construction. For the solution of the
extremization problem using Lagrangians and Lagrange multipliers see Mead
\etal\ \citey{meadetal1984} and the extensive texts by Fang \etal\
\citey{fangetal1997} and Boyd \etal\ \citey{boydetal2004_r2009}. For a
geometric understanding of the extremization and of the relation between
expectations and multipliers see Porta Mana \citey{portamana2017b}.

The result has the standard exponential-family form
\begin{equation}
  \label{eq:app_maxent_pop}
  \begin{gathered}
  \px(\yA)  = \frac{1}{Z(\yll)}\;
  r_{\yA}\;
  \exp\Biggl[\sum_{m}\yl_{m}\;
  \binom{\yA}{m}\binom{N}{m}^{-1}
  \Biggr],
  \\
  Z(\yll)\defd \sum_{\yA} r_{\yA}\;
  \exp\Biggl[\sum_{m}\yl_{m}\;
  \binom{\yA}{m}\binom{N}{m}^{-1}
  \Biggr],
\end{gathered}
\end{equation}
where $(r_{\yA})$ is the reference distribution and $\yll\defd (\yl_{m})$ are
the Lagrange multipliers, satisfying the implicit
constraint equations~\eqref{eq:final_constraint}:
\begin{equation}
  \label{eq:constraint_eq_lambda}
  \sum_{\yA}
  \tbinom{\yA}{m}\tbinom{N}{m}^{-1}\;
  \tfrac{1}{Z(\yll)}\;
  r_{\yA}\;
  \exp\biggl[\sum_{m}\yl_{m}\;
  \tbinom{\yA}{m}\tbinom{N}{m}^{-1}
  \biggr]
=
  \sum_{\ya}
  \tbinom{\ya}{m} \tbinom{n}{m}^{-1} \; \yff_{\ya},\quad
  m\in M.
\end{equation}

The reference distribution $(r_{\yA})$ represents our pre-data beliefs
about the activity levels $\yA$. We know that the majority of neurons in a
brain area are rarely simultaneously active within a window of some
milliseconds, so we could choose a distribution with slightly higher
weights on low values of $\yA$. On the other hand, considering the number
of ways in which $\yA$ out of $N$ neurons can be simultaneously active
would suggest the multiplicity distribution proportional to
$\binom{N}{\yA}$. It turns out that our results of
\sects~\ref{sec:application}--\ref{sec:marginalization} are insensitive to
the choice between these two possible reference distributions, or even a
uniform reference distribution.

For a derivation from a Bayesian perspective see
appendix~\ref{sec:inference_prob}.


\section{Measure of informational sufficiency}
\label{sec:measure_suff}

Let's ask how much more probable is the
sufficiency of one set with respect to the other, conditional on our data
$\yf$:
\begin{equation}
  \label{eq:ratio_prob_suffic}
  \frac{p(M'' \| \yf)}{ p(M' \| \yf)}.
\end{equation}
The probability of observing activity $\ya$ in the sample at any time bin
is the sample marginal of the maximum-entropy distribution for the larger
population, owing to the excheangeability assumption implicit in the
maximum-entropy method:
\begin{equation}
  \label{eq:sample_marginal}
  \pxx(\ya \| M) = \sum_{\yA}\yG_{\ya\,\yA}\;\px(\yA \| M).
\end{equation}
The probability of observing a specific sequence $(\ya_{t})$ with
frequencies $\yf$ is therefore
\begin{equation}
  \label{eq:prob_frequencies}
  \prod_{t=1}^{T} \pxx(\ya_{t} \| M) \equiv
  \prod_{\ya=0}^{n} \pxx(\ya \| M)^{T\,\yff_{\ya}} \equiv
  \prod_{\ya=0}^{n}\biggl[\sum_{\yA}\yG_{\ya\yA}\;\px(\yA \| M)\biggr]^{T\,\yff_{\ya}}.
\end{equation}
The probability of observing the frequencies $\yf$ is obtained multiplying this
by their multiplicity factor, the multinomial coefficient
\begin{equation}
  \label{eq:multiplicity_factor}
  \binom{T}{T\yf} \defd \frac{T!}{\prod_{\ya}(T\,\yff_{\ya})!}
  \approx \prod_{\ya}{\yff_{\ya}}^{-T\,\yff_{\ya}},
\end{equation}
the last expression coming from Stirling's approximation
\citep[Lemma~2.2]{csiszaretal2004b}.
If we assign equal pre-data probabilities to the two hypotheses $M'$ and $M''$,
each probability in the ratio~\eqref{eq:ratio_prob_suffic} then becomes, by
Bayes's theorem,
\begin{multline}
  \label{eq:prob_proportional_relentropy}
  \pf(M \| \yf) \propto
  \pf(\yf \| M) % \times\text{\small const}
  \propto
  \binom{T}{T\yf}\,  \prod_{\ya}\biggl[\sum_{\yA}\yG_{\ya\yA}\px(\yA \| M)\biggr]^{
    T\,\yff_{\ya}}
   \approx{}\\
  \prod_{\ya}{\yff_{\ya}}^{-T\,\yff_{\ya}}
  \times
  \prod_{\ya}\biggl[\sum_{\yA}\yG_{\ya\yA}\px(\yA \| M)\biggr]^{
    T\,\yff_{\ya}}.
\end{multline}
The logarithm of the probability above is easily seen to be the number of
bins $T$ multiplied by relative entropy between the frequency distribution
$\yf$ and the sample marginal of the maximum-entropy distribution.

Thus, the difference~\eqref{eq:measure_suffic} is the logarithm of the
probability ratio~\eqref{eq:ratio_prob_suffic}. The exponential of the
difference~\eqref{eq:measure_suffic} tells us how much more probable the
hypothesis that the set $M''$ is sufficient than the hypothesis that the
set $M'$ is.

\section{Assumptions and approximations: mathematical sketch}
\label{sec:inference_prob}

We give here a sketch of the mathematical formulae behind the assumptions
and approximations discussed in \sect~\ref{sec:assumptions}. We indicate
them all collectively with $\yI$.

The first assumption, of infinite exchangeability
\citep[\chap~4]{bernardoetal1994_r2000}{dawid2013%cifarellietal1982
} of the probability distribution $\pf[(\yA_{1},\dotsc, \yA_{T}) \| \yI]$
for the sequence of activities, implies the following mixture form by
de~Finetti's theorem \citep{definetti1930,hewittetal1955}:
\begin{equation}
  \label{eq:mixture_from_definetti}
%  \exp[-T\;\sh(\yF;\ynu)] \quad
%  \pxxx[(\yA_{1},\dotsc,\yA_{T})] \propto
  \pf[(\yA_{1},\dotsc, \yA_{T}) \| \yI]
  = \int\!\!\!\di\ynu\;\pf(\ynu \|\yI)\,\prod_{\yA=0}^{N}{\nu_{\yA}}^{T\,\yFF_{\yA}},
\end{equation}
where $\ynu\defd(\nu_{\yA})$ is the long-run frequency distribution of
activities in an imaginary continuation or repetition of the recording, and
$\pf(\ynu \|\yI)\;\di\ynu$ our belief distribution about it. The
integration is over an $N$-dimensional simplex.

The second assumption concerns the latter distribution. Its density
function has one of the following expressions:
\begin{enumerate}[label=(\textit{\alph*}),wide]
\item an entropic prior \citep{skillingetal1984,rodriguez1991,neumann2007}
  \begin{subequations}\label{eq:maxent_prior}
    \begin{equation}
      \label{eq:entropic_prior}
      \pf(\ynu \| \yI) \propto  \exp[-\yL\;\sh(\yF;\yr)]
      \approx \binom{L}{L\nu_{0}, \dotsc, L\nu_{N}}\;
      \prod_{\yA}{r_{\yA}}^{L\nu_{\yA}}  
    \end{equation}
    where $\yr\defd(r_{\yA})$ is a reference distribution and the
    expression in larger parentheses is a multinomial coefficient;
  \item a \enquote{sufficiency prior}, from the exponential family, which
    follows from the Pitman-Koopman theorem if some moments $\set{m}$ are
    sufficient statistics for our belief about $\ynu$:
    \begin{equation}
      \label{eq:suff_prior}
      \pf(\ynu \| \yI) \propto
      \int\!\!\!\di\yla\;\pf(\yla \| \yI)\;
      \prod_{\yA}\delt\Biggl\{\ynu_{\yA}-
      \frac{r_{\yA}}{Z(\yla)}\,\exp\Bigl[
      \tsum_{m}\lambda_{m}\tbinom{\yA}{m}/\tbinom{N}{m}
      \Bigr]\Biggr\},
    \end{equation}
  \end{subequations}
  where $\yla\defd(\lambda_{m})$ are Lagrange multipliers and $Z(\yla)$ a
  normalization factor for the exponential distribution within the delta.
  This is a formal expression, the delta symbolizing that the possible
  domain of $\ynu$ is restricted from an $N$-dimensional simplex to a
  lower-dimensional subset thereof, parametrized by $\yla$. When this prior
  is used, all integrations over $\ynu$ are replaced by integrations over
  $\yla$.
\end{enumerate}


The third assumption corresponds to the formula
\begin{multline}
  \label{eq:assume_sampled_anew}
  \pf[(\ya_{1},\dotsc,\ya_{T}) \| (\yA_{1}, \dotsc, \yA_{T}), \yI ]
  ={}\\ \prod_{t=1}^{T} \pf(\ya_{t} \| \yA_{t}, \yI)
  = \prod_{t=1}^{T} \yG_{\ya_{t}\yA_{t}}
  \equiv \prod_{\ya,\yA}{\yG_{\ya_{t}\yA_{t}}}^{T\,J_{\ya\yA}},
\end{multline}
where $J_{\ya\yA}$ is the joint frequency distribution of the activity
pairs $(\ya,\yA)$, so that its marginals are the frequency distributions
$\yf$ and $\yF$. As discussed in \sect~\ref{sec:assumptions} this formula
follows from either of two beliefs: that all subsets of $\yA$ neurons in
the larger population are always equally likely to give rise to total
activity $\yA$, or that a new sampling is done at every time bin.

It can be proved that the formulae above lead to the following belief
distribution for the joint frequency distribution $\yJ \defd (J_{\ya\yA})$
and the long-run frequency distribution $\ynu$, conditional on $\yf$:
\begin{multline}
  \label{eq:prob_joint_frequencies}
  \pf(\yJ, \ynu \| \yf, \yI) \propto
  %\int\!\!\!\di\ynu\;
  \pf(\ynu\|\yI)\;
                         \binom{T}{TJ_{0 0},\dotsc,TJ_{n N}}\,
                         \prod_{\ya,\yA}\bigl(\yG_{\ya\yA}\;\nu_{\yA}\bigr)^{T\,J_{\ya\yA}}
  \approx{}  \\
  \shoveright{%\int\!\!\!\di\ynu\;
    \pf(\ynu\|\yI)\;
                         \exp\biggl(
-T \tsum_{\ya,\yA}J_{\ya\yA}\ln\frac{J_{\ya\yA}}{\yG_{\ya\yA}\;\nu_{\yA}}
                         \biggr)}
  \\
  \text{with}\quad \tsum_{\yA}J_{\ya\yA}=\yff_{\ya},
\end{multline}
where $\pf(\ynu\|\yI)$ is either of the
distributions~\eqref{eq:maxent_prior}. Note the relative entropy in the
exponential. If we want to condition not on $\yf$ but on a set of measured
averages $C(m) = \sum_{\ya}c_{\ya}(m) \yff_{\ya}$, $m\in M$, such
as~\eqref{eq:chain_prescription_moments}, we have to replace the
integration constraint $\sum_{\yA}J_{\ya\yA}=\yff_{\ya}$ with
$\sum_{\ya\yA}c_{\ya}(m)\,J_{\ya\yA}=C(m)$.

From this belief distribution we can find the one for
$\yFF_{\yA}\equiv \sum_{\ya}J_{\ya\yA}$ and any others of interest by
marginalization. Within the exponential we can recognize the relative
entropy of $\yJ$ with respect to $(\yG_{\ya\yA}\;\nu_{\yA})$, which is the
reason why these two joint distributions must be equal if $T$ tends to
infinity.

From the formulae above we can also find the four distributions mentioned
at the end of \sect~\ref{sec:model}. The frequency
distribution~\ref{item:prob_freq} is the mode of $\pf(\yF \| \yf,\yI)$. The
belief distribution~\ref{item:prob_binin} is
\begin{equation}
  \label{eq:prob_binin}
  \pf(\yA_{t} \| \yf, \yI) =
  \sum_{\yF} \yFF_{\yA_{t}}\; \pf(\yF \| \yf, \yI),
  \qquad t \in \set{1,\dotsc,T}.
\end{equation}
The frequency distribution~\ref{item:prob_longrun} is the mode of $\pf(\ynu
\| \yf,\yI)$. Finally, the belief distribution~\ref{item:prob_binout} is
\begin{equation}
  \label{eq:prob_binout}
  \pf(\yA_{t} \| \yf, \yI) =
  \int\!\!\!\di\ynu\;\nu_{\yA_{t}}\;\pf(\ynu \| \yf, \yI),
  \qquad t \notin \set{1,\dotsc,T}.
\end{equation}




%%%%%%%%%%%%%%%%%%%%%%%%%%%%%%%%%%%%%%%%%%%%%%%%%%%%%%%%%%%%%%%%%%%%%%%%%%%%
%%% Bibliography
%%%%%%%%%%%%%%%%%%%%%%%%%%%%%%%%%%%%%%%%%%%%%%%%%%%%%%%%%%%%%%%%%%%%%%%%%%%% 
\clearpage
\defbibnote{prenote}{{\footnotesize (\enquote{de $X$} is listed under D,
    \enquote{van $X$} under V, and so on, regardless of national
    conventions.)\par}}
% \defbibnote{postnote}{\par\medskip\noindent{\footnotesize% Note:
%     \arxivp \mparcp \philscip \biorxivp}}

\printbibliography[prenote=prenote%,postnote=postnote
]

\end{document}

%%%%%%%%%%%%%%%%%%%%%%%%%%%%%%%%%%%%%%%%%%%%%%%%%%%%%%%%%%%%%%%%%%%%%%%%%%%%
%%% Cut text (won't be compiled)
%%%%%%%%%%%%%%%%%%%%%%%%%%%%%%%%%%%%%%%%%%%%%%%%%%%%%%%%%%%%%%%%%%%%%%%%%%%% 

%%% Local Variables: 
%%% mode: LaTeX
%%% TeX-PDF-mode: t
%%% TeX-master: t
%%% End: 
